\documentclass[11pt,letterpaper]{article}
\usepackage[utf8]{inputenc}
\usepackage[T1]{fontenc}
\usepackage[spanish,es-noshorthands]{babel}
\usepackage[margin=1in]{geometry}
\usepackage{times}
\usepackage{microtype}
\usepackage{graphicx}
\usepackage{xcolor}
\usepackage{booktabs}
\usepackage{longtable}
\usepackage{array}
\usepackage{hyperref}
\hypersetup{colorlinks=true, urlcolor=blue, linkcolor=blue, citecolor=blue}
\usepackage{listings}
\usepackage{titlesec}
\usepackage{enumitem}
\usepackage{caption}
\captionsetup{font=small}
\usepackage{tikz}
\usetikzlibrary{shapes.geometric,arrows.meta,positioning,decorations.pathreplacing,backgrounds,fit,calc,patterns}
\usepackage{pgfplots}
\pgfplotsset{compat=1.18}
\usepackage{setspace}
\linespread{1.10}

% Typography / page-break behaviour
\raggedbottom
\sloppy
\widowpenalty=10000
\clubpenalty=10000

% Top-level \section{...} should start on a fresh page
\let\oldsection\section
\renewcommand{\section}{\clearpage\oldsection}

\lstset{
  basicstyle=\ttfamily\footnotesize,
  breaklines=true,
  frame=single,
  showstringspaces=false,
  numbers=left,
  numberstyle=\tiny\color{gray},
  numbersep=6pt,
  xleftmargin=2em
}

\tikzset{
  cdmbox/.style={draw, rounded corners=2pt, minimum height=0.8cm, font=\small, align=center},
  thinarrow/.style={-{Stealth[length=2mm,width=1.5mm]}, thick},
  fatarrow/.style={-{Stealth[length=3mm,width=2mm]}, very thick},
  layerbox/.style={draw, rounded corners=3pt, minimum height=1.4cm, minimum width=9cm, font=\small, align=left, inner sep=4pt},
  slotcell/.style={draw, minimum width=0.85cm, minimum height=0.7cm, font=\scriptsize, align=center, inner sep=2pt},
}

\title{Adentro de Widevine L1 en Windows:\\
El despacho de tres capas que nadie mapeó,\\
y las ranuras por-cuadro que nunca disparan}

\author{
  Jeremy Erazo (\texttt{trexnegr0}) \\
  Investigador independiente de seguridad \\
  \texttt{mendozayt13@gmail.com}
}

\date{Junio de 2026}

\begin{document}
\maketitle

\vspace{-10pt}
\begin{center}
\fbox{\begin{minipage}{0.92\columnwidth}\small
\textbf{Resumen ejecutivo.}
El CDM Widevine L1 que viene en todo navegador de la familia
Chromium en Windows ha sido, hasta ahora, una caja negra para
cualquiera fuera de Google. Este paper la abre. Tres vtables
concéntricas, diecinueve ranuras públicas de CDM\_11, verificadas
byte-a-byte en dos builds de navegadores distintos. Un mapa
empírico runtime-confirmado de qué ranuras se ejercen realmente
durante streaming hardware-decoded real. Una distribución de
ofuscación estilo MBA a lo largo del espacio de ranuras que
refleja la profundidad del parsing de datos atacante-controlables.
Tablas de alcanzabilidad tanto para Netflix como para una fuente
pública Widevine no-Netflix. Y, de fondo, una metodología que
captura todo esto sin jamás leer un byte de contenido, una clave de
contenido, ni un IV desde la memoria del CDM. Los defensores se
llevan una tabla de superficie de ataque por ranura más firmas de
detección; los atacantes se llevan la desambiguación de cuáles
ranuras tienen que vivirles. La metodología completa se libera con
fuente junto al paper; no se libera artefacto de extracción ni
exploit weaponizado.
\end{minipage}}
\end{center}
\vspace{8pt}

\begin{abstract}
\noindent
El Módulo de Descifrado de Contenido (CDM) Widevine viene
prácticamente en todo navegador de escritorio mayor y soporta el
camino de reproducción de Encrypted Media Extensions (EME) para
todo servicio comercial de streaming. Los internals del build L1
de Windows están documentados sólo fragmentariamente: los archivos
de cabecera describen el ABI de C++ que expone al navegador
anfitrión, pero el binario en sí es código cerrado,
parcialmente ofuscado, y se comporta diferente dependiendo de si
el anfitrión delega la decodificación a pipelines protegidos por
hardware. Este paper presenta una ingeniería inversa arquitectural
estilo Buchanan de la implementación CDM\_11 tal como viene en
Chrome 149 y Edge 150 sobre Windows 11. Identificamos un
\emph{modelo de despacho de tres capas} (\emph{thunks} de la
vtable externa, una vtable intermedia \texttt{m\_impl} de
\emph{wrappers} limpios y \emph{double-trampolines}, y cuerpos de
\emph{handlers} por ranura), localizamos cada capa byte-a-byte en
ambos binarios, y caracterizamos empíricamente cuáles ranuras se
invocan durante reproducción en steady-state vs. cuáles permanecen
dormidas. Confirmamos en runtime que durante reproducción
hardware-decoded de Netflix las ranuras por-cuadro del CDM\_11
\texttt{Decrypt}, \texttt{DecryptAndDecodeFrame},
\texttt{DecryptAndDecodeSamples}) no disparan mientras una tabla
de despacho singleton interna genera $\sim$2{,}300 llamadas por
segundo de trabajo cripto/scheduler invisible desde el límite
EME. Documentamos la distribución de ofuscación estilo
mixed-boolean-arithmetic en el espacio de ranuras y mostramos que
se concentra donde el CDM empieza a parsear datos estructurados
atacante-controlables, mientras está ausente en ranuras que pasan
\emph{blobs} opacos. Después verificamos la alcanzabilidad de las
ranuras de gestión-de-sesión desde dos aplicaciones EME distintas
(Netflix y una fuente pública de prueba Widevine) y mostramos
comportamiento de despacho idéntico. El paper sitúa este trabajo
en dos décadas de investigación previa de DRM, recorre los
precedentes legales y metodológicos, y documenta tanto los caminos
exitosos como los falsos arranques que caracterizaron el proceso de
ingeniería inversa. Ningún clave de contenido, buffer de cuadro, ni
byte de ciphertext se captura en ningún momento del estudio.
\end{abstract}

% =============================================================================
\section{Introducción}
\label{sec:intro}

Modern desktop video DRM is the bridge between commercial streaming
services and the browser's media pipeline. When a user plays a video
in Netflix, Amazon Prime, Disney+ or any other Widevine-protected
source, the browser invokes Encrypted Media Extensions (EME) to
delegate the cryptographic operations to a Content Decryption Module
(CDM). In Chromium-based browsers on Windows, that CDM is the
proprietary \texttt{widevinecdm.dll} library shipped by
Google~\cite{widevine-spec}. The browser loads the library inside a
Mojo utility process
(\texttt{--utility-sub-type=media.mojom.\\CdmServiceBroker}), connects
to it over IPC, and asks the library to construct a CDM instance
implementing a stable C++ interface defined in
\texttt{content\_decryption\_module.h}~\cite{chromium-cdm-api}.
Desde el punto de vista del anfitrión, la interfaz es chica
(diecinueve métodos virtuales) y bien documentada:
\texttt{Initialize}, \texttt{GetStatusForPolicy},
\texttt{CreateSessionAndGenerateRequest},
\texttt{UpdateSession}, y otros doce métodos que cubren el
intercambio de licencias, el descifrado, la decodificación y el
lifecycle.

Desde la perspectiva de la \emph{biblioteca}, esta interfaz limpia
de diecinueve métodos es la punta visible de un iceberg mucho
mayor. El build L1 de Windows de Widevine implementa un runtime
ofuscado de varios millones de instrucciones bajo decenas de
megabytes de código cerrado, se integra con Windows Protected
Media Path (PMP) cuando hay decodificación por hardware
disponible, y está sujeto a un régimen de licencia estricto que
prohíbe el estudio de sus internals de descifrado. El conocimiento
público sobre la implementación es fragmentario: el árbol de
Chromium contiene el \emph{lado anfitrión} del cableado EME,
varios proyectos académicos y de aficionados han estudiado la
variante \emph{Android L3} \cite{quarkslab-widevine,
l3-decryptor}, pero la variante L1 de escritorio en Windows ha
recibido comparativamente poca atención arquitectural en la
literatura abierta.

La motivación de este estudio es tanto metodológica como técnica.
Varios episodios legales recientes han mostrado que
\emph{publicar lo que uno aprende} sobre internals de DRM está más
expuesto legalmente que \emph{aprenderlo}, y que la
reproducibilidad de un resultado pesa tanto como el resultado en
sí \cite{sklyarov-case, dvdjon-norway, hadad-case}. Siguiendo el
modelo Buchanan de ingeniería inversa arquitectural
\cite{buchanan-blog, buchanan-imessage} --- explorar todo,
documentar la técnica, no publicar artefacto de extracción ---
nos propusimos construir una metodología completamente
reproducible y completamente abierta para estudiar el CDM
Widevine y aplicarla a las preguntas que son publicables sin
cruzar la línea de la circumvención de gestión de derechos
digitales.

\subsection*{Lo nuevo de este trabajo}

\begin{enumerate}[itemsep=4pt,topsep=2pt]
  \item \textbf{El primer mapa arquitectural público de L1
        \texttt{widevinecdm.dll} en Windows.} Tres vtables
        concéntricas, cada entrada clasificada, cada offset de
        despacho verificado byte-a-byte. Validación cruzada en
        Chrome 149 y Edge 150. Reproducible drop-in sobre cualquier
        build futuro vía el discoverer libre de parámetros que
        liberamos.

  \item \textbf{Inactividad bajo decode por hardware,
        runtime-confirmada.} Cuatro ranuras de la outer-vtable
        patcheadas en un solo experimento. A lo largo de más de 90
        segundos de reproducción 4K HDR continua de Netflix,
        \(g\_calls = 0\) en la ranura 14
        (\texttt{DecryptAndDecodeFrame}), la 15
        (\texttt{DecryptAndDecodeSamples}) y la 9
        (\texttt{Decrypt}) simultáneamente. El trabajo por-cuadro
        en L1 no atraviesa la interfaz pública del CDM; la
        superficie de ataque vista desde EME para intercepción
        por-cuadro es, en términos prácticos, vacía.

  \item \textbf{Una tabla de superficie de ataque por ranura.}
        Qué ranuras ingieren input estructurado
        atacante-controlable (PSSH, respuesta de licencia,
        política), qué ranuras pasan blobs opacos, qué ranuras
        están dormidas. Consistente entre fuentes: Netflix y el
        demo Widevine de Bitmovin ejercen las mismas ranuras en el
        mismo orden. Los defensores usan la tabla directamente
        para alcance de detección; los atacantes la usan
        directamente para priorizar los handlers de capa-3 que de
        verdad importan.

  \item \textbf{Ofuscación estilo MBA, mapeada relativa a la
        profundidad.} Cada una de las diecinueve ranuras CDM\_11
        caracterizada en cada capa del despacho. La ofuscación
        vive a profundidades variables, apareciendo consistentemente
        donde los datos estructurados del atacante empiezan a ser
        parseados y ausente consistentemente en ranuras de
        passthrough. Ejemplos concretos con firmas byte-level para
        cada una.

  \item \textbf{Una metodología que viaja.} El patrón de
        \emph{observación shape-only} es generalizable a cualquier
        binario closed-source con una interfaz estable. El mismo
        diseño de hook aplica, sin cambiar de forma, al
        \texttt{PlayReadyHost} de PlayReady, al
        \texttt{AppleSpringboard} de FairPlay, o a cualquier otro
        contrato host-CDM donde la pregunta es \emph{qué entradas
        disparan cuándo} y la respuesta no debe capturar
        contenido. Liberamos el framework estilo Buchanan como una
        contribución metodológica en sí misma.

  \item \textbf{Manual del defensor y lectura del atacante.} Dos
        secciones explícitas que traducen los hallazgos en acción
        concreta. Firmas de detección, canarios de hook, perfiles
        de comportamiento, y priorización por ranura para ambos
        lados.

  \item \textbf{Reproducible y falsable.} Cada afirmación empírica
        está anclada a una corrida experimental grabada. Cada
        afirmación arquitectural es rederivable desde el binario
        público mediante el discoverer que liberamos. El sprint de
        falsificación que precedió a la publicación está
        documentado en su totalidad.
\end{enumerate}

\subsection*{Hoja de ruta del paper}

El paper procede de la siguiente forma. La
Sección~\ref{sec:antecedents} sitúa este trabajo dentro de dos
décadas de ingeniería inversa previa de DRM comercial y los
precedentes legales que la han moldeado.
La Sección~\ref{sec:background} introduce los niveles de licencia
de Widevine, el camino de despacho EME, la arquitectura Mojo IPC
que usa Chromium para hospedar el CDM, el límite PMP de Windows,
el formato Common Encryption (CENC) que usa todo el contenido
Widevine, y la estructura de la caja PSSH que transporta datos
de inicialización atacante-controlables.
La Sección~\ref{sec:threat} declara el modelo de amenaza y el
alcance ético.
La Sección~\ref{sec:static} recorre la ingeniería inversa estática
de CDM\_11 en detalle, con listings de disassembly concretos.
La Sección~\ref{sec:runtime} describe la instrumentación en
runtime, incluyendo las sutilezas de diseño del bypass de CIG, el
registro de bitmap CFG, y la cadena de back-pointers por ranura.
La Sección~\ref{sec:findings} presenta los hallazgos empíricos.
La Sección~\ref{sec:obfuscation} trata la distribución de
ofuscación como su propia defensa arquitectural.
La Sección~\ref{sec:cross} establece la estabilidad
inter-binarios e inter-fuentes.
La Sección~\ref{sec:robustness} reporta sobre las pruebas de
falsificación, incluyendo los hallazgos candidatos que fueron
refutados por análisis más profundo.
La Sección~\ref{sec:discussion} discute lo que el modelo implica
tanto para atacantes como defensores del DRM de escritorio.
La Sección~\ref{sec:ethics} declara qué no se hizo y por qué.
La Sección~\ref{sec:conclusion} concluye.

\begin{figure}[t]
\centering
\begin{tikzpicture}[node distance=0.6cm and 1.2cm]
  \node[cdmbox, fill=blue!10, minimum width=3.2cm] (renderer) {Renderer process\\\scriptsize JavaScript + EME};
  \node[cdmbox, fill=gray!15, right=of renderer, minimum width=3.2cm] (browser) {Browser process\\\scriptsize CDM dispatcher};
  \node[cdmbox, fill=orange!15, below=of browser, minimum width=4.2cm] (utility) {CDM utility process\\\scriptsize \texttt{--utility-sub-type=}\\\scriptsize \texttt{media.mojom.CdmServiceBroker}};
  \node[cdmbox, fill=red!12, below=of utility, minimum width=4.2cm] (wvdll) {\texttt{widevinecdm.dll}\\\scriptsize 19-method CDM\_11 interface};
  \node[cdmbox, fill=yellow!20, below=of wvdll, minimum width=4.2cm] (cdmobj) {CDM\_11 instance\\\scriptsize 3-layer dispatch (this paper)};

  \draw[thinarrow] (renderer) -- node[above,font=\scriptsize] {EME events} (browser);
  \draw[thinarrow] (browser.south) -- node[right,font=\scriptsize,xshift=2pt] {Mojo IPC} (utility.north);
  \draw[thinarrow] (utility) -- node[right,font=\scriptsize,xshift=2pt] {\texttt{LoadLibrary}} (wvdll);
  \draw[thinarrow] (wvdll) -- node[right,font=\scriptsize,xshift=2pt] {\texttt{CreateCdmInstance(11)}} (cdmobj);
\end{tikzpicture}
\caption{Despacho Widevine EME extremo-a-extremo en navegadores
de la familia Chromium sobre Windows. El navegador anfitrión
delega las operaciones de Encrypted Media Extensions vía Mojo IPC
a un proceso utility cuyo único propósito es cargar
\texttt{widevinecdm.dll} y exponer su interfaz CDM\_11. Este paper
estudia el modelo de despacho que se ubica entre la interfaz
devuelta por \texttt{CreateCdmInstance} y los handlers por-ranura
dentro de la biblioteca.}
\label{fig:system-overview}
\end{figure}

% =============================================================================
\section{Antecedentes y trabajos relacionados}
\label{sec:antecedents}

La ingeniería inversa de DRM ocupa un rincón inusual de la
investigación en seguridad: abarca criptografía aplicada,
ingeniería inversa de binarios, derecho de copyright, relaciones
con vendors, y (a veces) prosecución penal. La metodología que
seguimos en este paper no surgió en el vacío; está moldeada por
dos décadas de trabajo previo, parte del cual produjo resultados
técnicos excelentes y parte del cual produjo resultados técnicos
excelentes \emph{y} consecuencias legales adversas para el
investigador. Esta sección recorre el estado del arte previo y
explica qué precedentes informaron las decisiones metodológicas
tomadas aquí.

\subsection{El modelo Buchanan}
\label{sec:ante-buchanan}

El marco metodológico que seguimos está más claramente
articulado en el trabajo reciente de David Buchanan
\cite{buchanan-blog, buchanan-imessage, buchanan-t2}. La
investigación de Buchanan incluye notablemente la ingeniería
inversa del envoltorio de clave basado en NaCl de iMessage de
Apple, la cadena de arranque seguro del chip de seguridad T2 de
Apple, y un bug arbitrario en imágenes en la cadena del decoder
de imágenes de iOS (CVE-2023-4863).

La propiedad definitoria del trabajo publicado de Buchanan, para
nuestros fines, es la separación estricta entre
\emph{descripción arquitectural} y \emph{artefacto explotable}.
Su output público documenta cómo está estructurado un sistema y
dónde están sus fronteras de confianza; las pruebas de concepto
controladas que acompañan los writeups demuestran la clase de bug
sin producir una herramienta que directamente habilite el abuso
contra víctimas no relacionadas. Él describió públicamente la
disciplina de publicación resultante como: explorar todo,
documentar la técnica, publicar la arquitectura, no publicar el
artefacto que permitiría a cualquier lector materializar el
impacto.

Para investigación de DRM esta distinción importa: el 17~U.S.C.~§1201
(la provisión anti-circumvención del DMCA) y sus equivalentes
internacionales prohíben la publicación de herramientas o
tecnologías cuyo uso \emph{principal} sea la circumvención de una
medida tecnológica de protección. Una descripción arquitectural
de un CDM no es esa clase de herramienta. Un frame ripper que
consume la misma descripción arquitectural y produce MP4s de
Netflix descifrados claramente sí lo es. El espacio entre ambos
es la línea que cuidadosamente no cruzamos.

\subsection{Hacking the Xbox y la tradición bunnie}
\label{sec:ante-bunnie}

El trabajo de Andrew ``bunnie'' Huang sobre ingeniería inversa de
consolas --- la Xbox original \cite{bunnie-xbox-book}, el pipeline
de video con manejo de HDCP del proyecto NeTV
\cite{bunnie-nettv}, y la laptop Novena --- demostró un modelo
sustentable de investigación adyacente al DRM que era
simultáneamente útil, publicable, y no objeto de litigación por
parte del vendor. El proyecto NeTV en particular es significativo:
HDCP es una medida de protección de copyright bajo la ley
estadounidense, y el proyecto requirió correspondencia coordinada
con la EFF dirigida a la Library of Congress para confirmar que
una demostración arquitectural que no captura, redistribuye, ni
facilita la redistribución de contenido protegido se mantiene
legal bajo el marco de excepciones del DMCA.

La observación recurrente de Huang es que la exposición legal se
adhiere a la \emph{exhibición} de una capacidad de circumvención
y a la \emph{distribución} de herramientas que la realizan, no al
acto de caracterizar un sistema. El modelo de despacho CDM\_11
descrito en este paper sería obviamente útil para un atacante que
ya poseyera una capacidad independiente de recibir contenido
Widevine cifrado y obtener licencias válidas; sería esencialmente
inútil para un atacante casual que no la tuviera. Esa distribución
de utilidad está, en nuestra opinión, del lado correcto de la
línea bunnie.

\subsection{Sklyarov, DVD-Jon, y la frontera legal}
\label{sec:ante-legal}

Los casos de \textbf{Dimitry Sklyarov} y \textbf{Jon Lech
Johansen} delimitan la frontera legal de la investigación de DRM.

Sklyarov, un desarrollador ruso de ElcomSoft, fue arrestado en
2001 durante DEF~CON por su trabajo en el Advanced eBook
Processor, que podía convertir los eBooks protegidos de Adobe en
PDF~\cite{sklyarov-case}. La prosecución procedió contra
ElcomSoft (Sklyarov fue eventualmente liberado tras aceptar
testificar). El arresto, el juicio, y la eventual absolución de
ElcomSoft sobre los cargos penales establecieron dos precedentes
relevantes: (i) los fiscales van a tratar a una herramienta cuya
salida es la forma en texto plano de una obra protegida como
``usada primariamente para circunvenir'', aunque la investigación
subyacente sea académicamente defendible; (ii) meramente
\emph{describir} la técnica de circumvención, sin distribuir una
herramienta que la realice, no fue la base de la prosecución.

Jon Lech Johansen, ``DVD-Jon'', fue procesado por las autoridades
noruegas (absuelto dos veces, finalmente definitivamente) por su
rol en el desarrollo de DeCSS, la herramienta que rompió el
scrambling CSS en los DVDs~\cite{dvdjon-norway}. Los tribunales
noruegos sostuvieron que el trabajo de Johansen sobre su
\emph{propio} reproductor de DVD era legal, aunque la herramienta
resultante haya sido redistribuida ampliamente en internet. El
caso estadounidense de \emph{Universal v.~Reimerdes} sostuvo que
la \emph{publicación} de DeCSS en un sitio web era circumvención
accionable del DMCA, sin importar la legalidad de la
investigación subyacente~\cite{decss-2600}.

La lección compuesta de estos casos informa nuestra disciplina de
publicación: una descripción arquitectural no es una herramienta;
una herramienta que materializa contenido descifrado sí es una
herramienta. Publicamos la descripción arquitectural y el marco
de validación controlada que prueba que la descripción es
correcta; no publicamos una herramienta.

\subsection{Trabajo académico sobre Widevine}
\label{sec:ante-widevine}

El ecosistema Widevine ha recibido atención desde varios grupos
académicos y corporativos, principalmente enfocados en la
variante Android L3.

\textbf{Quarkslab} publicó una serie de análisis de Widevine en
Android beginning in 2018~\cite{quarkslab-widevine,
quarkslab-android-widevine}. Their work characterised the
OEMCrypto-protected components of the Widevine SDK and demonstrated
that the L3 keybox-based protection model relied substantially on
software obfuscation that could be peeled back with whitebox-crypto
técnicas de ataque. Los proyectos del L3 decryptor que siguieron
\cite{l3-decryptor, ango-l3-decryptor} --- la mayoría publicados
como open-source de calidad investigación --- demostraron
extracción extremo-a-extremo de contenido Widevine~L3 desde la
biblioteca protegida. La respuesta de Google fue deprecar el
soporte de L3 en navegadores web y acelerar la migración a L1.

\textbf{Check Point Research} ha analizado periódicamente Widevine
sobre handsets Android, identificando vulnerabilidades específicas
de extracción de keybox y clases de bugs en la TrustZone de
OEMCrypto~\cite{checkpoint-widevine}. Sus reportes fueron
coordinados con Google y el vendor antes de su publicación.

\textbf{Romain Thomas y colaboradores} de Quarkslab han publicado
extensamente sobre herramientas prácticas de ingeniería inversa
que usamos aquí, incluyendo la biblioteca LIEF para parsing
multiplataforma de binarios \cite{lief-tool} y varios posts de
blog sobre whitebox crypto en DRM comercial. Su trabajo en
herramientas está directamente upstream de la metodología usada
en este paper para análisis estático de \texttt{widevinecdm.dll}.

No existe, hasta donde sabemos, ningún estudio académico publicado
sobre los internals de despacho del CDM L1 de Windows como tal. El
trabajo más cercano en espíritu es un puñado de posts de blog y
threads de foros que documentan issues específicos (por ejemplo,
el layout de símbolos de \texttt{libwidevinecdm.so} en Linux, los
cambios de ABI del CDM entre versiones mayores) sin proponer un
modelo arquitectural generalizable. Esa es la brecha que el
presente paper intenta llenar para el build L1 de Windows.

\subsection{Criptografía whitebox y ofuscación MBA}
\label{sec:ante-whitebox}

Los patrones de ofuscación observados en los handlers de capa-3
de CDM\_11 --- constantes mágicas de alta entropía escritas al
stack en la entrada de la función, instrucciones de avance de
estado por hash multiplicativo, cascadas \texttt{cmovae} sin
branches, rellenos de poison \texttt{0xaaaaaaaa} --- coinciden con
la firma de libro de texto de la ofuscación
\emph{mixed-boolean-arithmetic} (MBA) tal como se trata en la
literatura académica~\cite{eyrolles-mba, biondi-mba}. La ofuscación
MBA transforma una operación aritmética en una expresión
equivalente que mezcla operaciones bit-a-bit y aritméticas de modo
que derrota la simplificación algebraica sobre la forma ofuscada.
Se usa ampliamente en productos comerciales de protección de
software (VMProtect, Themida, Arxan) y en ofuscadores a medida
construidos por vendors para casos defensivos específicos.

Una técnica distinta pero relacionada, la \emph{criptografía
whitebox}, embebe material de clave criptográfica en
representaciones basadas en tablas que resisten la extracción
directa incluso cuando el atacante tiene acceso completo de
lectura al código ofuscado~\cite{chow-aes-whitebox,
lepoint-whitebox}. Los productos comerciales de DRM usan variantes
de whitebox-AES extensivamente; las claves cifradas en una
licencia Widevine L3 son desenvueltas por exactamente esa clase
de rutina antes de ser aplicadas al descifrado por-cuadro.

El patrón que reportamos en la Sección~\ref{sec:obfuscation} ---
ofuscación concentrada a la profundidad donde los datos
estructurados del atacante empiezan a ser parseados --- es
consistente con la prescripción de diseño de la literatura
whitebox: proteger al procesador criptográfico, no el pegamento
de despacho a su alrededor. Si la protección \emph{tiene éxito}
contra un atacante determinado es una pregunta separada que no
abordamos.

\subsection{Antecedentes de herramientas}
\label{sec:ante-tools}

Tres piezas concretas de herramientas en las que se apoya este
trabajo merecen reconocimiento:

\textbf{Capstone}~\cite{capstone-tool}, el framework de
disassembly multi-arquitectura usado por cada script en este
release que recorre la sección \texttt{.text} de
\texttt{widevinecdm.dll}.

\textbf{pefile}~\cite{pefile-tool}, el parser PE puro de Python
usado para enumerar exports, recorrer la tabla de secciones, y
leer RVAs desde \texttt{.rdata}.

\textbf{mingw-w64}, la cross-toolchain de Windows usada para
construir el shim y los helpers desde un host Linux. Notablemente,
el backend de gcc de mingw-w64 no soporta el structured exception
handling \texttt{\_\_try / \_\_except} de Microsoft, lo que nos
forzó a usar checks de seguridad basados en \texttt{VirtualQuery}
para los caminos de lectura del shim que usan registers como
target, en lugar de \texttt{memcpy} protegido por SEH. Esta es la
clase de pequeño detalle de implementación que no aparece en la
literatura llamativa de investigación de DRM pero que determina
si la metodología corre o no.

% =============================================================================
\section{Marco técnico}
\label{sec:background}

Esta sección revisa los componentes técnicos principales del
pipeline de reproducción Widevine L1 de Windows en mayor
profundidad de la que permitió la introducción.

\subsection{Niveles de licencia de Widevine}
\label{sec:bg-levels}

Widevine define tres niveles de seguridad anidados.

\textbf{L3} es un CDM puramente software. El estado criptográfico
se mantiene en memoria del proceso dentro de un runtime ofuscado.
Las claves de contenido se desenvuelven mediante un núcleo
criptográfico estilo whitebox y se aplican a segmentos de
ciphertext en buffers user-space. La salida descifrada se devuelve
al anfitrión a través de memoria estándar. Según la política de
licencia Widevine publicada, los clientes solo-L3 tienen techo en
540p; los acuerdos de los dueños de contenido frecuentemente
restringen L3 a tasas de bits más bajas que las que reciben los
clientes L2/L1.

\textbf{L2} es híbrido: un CDM software mantiene al menos sus
operaciones sensibles detrás de un Trusted Execution Environment
(TEE) respaldado por hardware en el anfitrión. L2 ve un deployment
limitado en escritorio.

\textbf{L1}, la variante estudiada en este paper, requiere que el
anfitrión provea un pipeline de medios protegido respaldado por
hardware. En Windows el pipeline relevante es el Protected Media
Path (PMP), él mismo un compuesto de:

\begin{itemize}[itemsep=0pt,topsep=2pt]
  \item \texttt{peauth.sys}, the Microsoft-signed kernel-mode
        component that enforces the OPM contract between the CDM
        and the GPU driver,
  \item the Media Foundation \texttt{IDirectXVideoDecoder} chain
        running in a kernel-protected mode,
  \item the DXGI protected-surface allocator that exposes the
        decoded frame only to the GPU presentation pipeline and not
        to user-space readback.
\end{itemize}

Under L1 hardware decode, the decrypted frame never appears in
user-mode browser process memory. This is the design property that
the paper's first empirical finding (Section~\ref{sec:findings-dormancy})
exploits to demonstrate the CDM dormancy phenomenon.

\subsection{Encrypted Media Extensions and the dispatch path}
\label{sec:bg-eme}

La especificación W3C Encrypted Media Extensions~\cite{w3c-eme}
define la API de JavaScript por la cual una aplicación web
solicita reproducción mediada por DRM de un stream de medios
cifrado. Desde el punto de vista del navegador anfitrión, las
operaciones relevantes son:

\begin{enumerate}[itemsep=2pt,topsep=2pt]
  \item \texttt{navigator.requestMediaKeySystemAccess(}
        \texttt{``com.widevine.alpha''}\texttt{)}: confirma que
        el cliente tiene un CDM para el key system nombrado.
  \item Crear un objeto \texttt{MediaKeys} y vincularlo a un
        \texttt{HTMLMediaElement}.
  \item Esperar el evento \texttt{encrypted}, que provee los
        datos de inicialización Common Encryption (CENC)
        extraídos del stream de medios.
  \item Crear una \texttt{MediaKeySession} y llamar a
        \texttt{generateRequest(initDataType, initData)}, lo que
        causa que el CDM construya un request de licencia.
  \item Reenviar el mensaje resultante al servidor de licencias,
        recibir la respuesta, y devolverla vía
        \texttt{updateSession(licenseResponse)}.
  \item Continuar la reproducción mientras el CDM desenvuelve
        las claves de contenido y las aplica a los datos del
        medio cifrado.
\end{enumerate}

La implementación de Chromium de esta API enruta cada paso a
través del proceso navegador, el renderer, y el proceso utility
del CDM vía interfaces Mojo IPC definidas en
\texttt{chromium/components/cdm}~\cite{chromium-cdm-mojo}. La
interfaz Mojo relevante para nuestros fines es
\texttt{media.mojom.ContentDecryptionModule}, cuyos mensajes de
request y response son esencialmente una proyección serializada
de la interfaz C++ de diecinueve métodos de CDM\_11.

\subsection{La interfaz C++ de CDM\_11}
\label{sec:bg-cdm11}

La interfaz que CDM\_11 implementa está definida en el árbol de
Chromium en
\texttt{media/cdm/api/content\_decryption\_module.h}. La clase es
una base puramente virtual \texttt{ContentDecryptionModule\_11}
con las siguientes diecinueve entradas, declaradas en este orden:

\begin{lstlisting}[basicstyle=\ttfamily\scriptsize]
virtual void Initialize(...) = 0;
virtual void GetStatusForPolicy(...) = 0;
virtual void SetServerCertificate(...) = 0;
virtual void CreateSessionAndGenerateRequest(...) = 0;
virtual void LoadSession(...) = 0;
virtual void UpdateSession(...) = 0;
virtual void CloseSession(...) = 0;
virtual void RemoveSession(...) = 0;
virtual void TimerExpired(...) = 0;
virtual Status Decrypt(...) = 0;
virtual Status InitializeAudioDecoder(...) = 0;
virtual Status InitializeVideoDecoder(...) = 0;
virtual void DeinitializeDecoder(...) = 0;
virtual void ResetDecoder(...) = 0;
virtual Status DecryptAndDecodeFrame(...) = 0;
virtual Status DecryptAndDecodeSamples(...) = 0;
virtual void OnPlatformChallengeResponse(...) = 0;
virtual void OnQueryOutputProtectionStatus(...) = 0;
virtual void OnStorageId(...) = 0;
\end{lstlisting}

Cada llamada virtual se traduce, bajo la convención de llamada
x64 de Microsoft, en una única llamada indirecta a través de un
function pointer cargado desde la vtable del objeto. El layout de
la vtable para una clase con diecinueve funciones virtuales y sin
virtuales de clase base es un arreglo contiguo de diecinueve
punteros a código; el overhead por llamada desde la perspectiva
del anfitrión es una carga indirecta más el costo del dispatcher
CFG.

Adoptamos la convención de identificar los métodos CDM\_11 por su
índice de ranura a lo largo de este paper. \texttt{Decrypt} es la
ranura 9, \texttt{DecryptAndDecodeFrame} es la ranura 14, y así.

\subsection{Common Encryption (CENC) y PSSH}
\label{sec:bg-cenc}

El protocolo Widevine transporta sus datos de inicialización
dentro de una caja Protection-Specific System Header (PSSH), tal
como está estandarizada en ISO/IEC 23001-7 (Common Encryption
para ISO Base Media File Format)~\cite{iso-cenc, w3c-pssh}. Una
caja PSSH está estructurada así:

\begin{lstlisting}[basicstyle=\ttfamily\scriptsize]
PSSHBox {
    u32         size;              // big-endian
    char[4]     type;              // "pssh"
    u8          version;
    u24         flags;
    UUID(16B)   SystemID;          // Widevine SystemID
                                   //  edef8ba9-79d6-4ace-a3c8-27dcd51d21ed
    if version > 0:
        u32     KID_count;
        UUID[KID_count] KIDs;
    u32         data_size;
    u8[data_size] data;
}
\end{lstlisting}

El campo \texttt{data}, para Widevine, contiene una serialización
Protocol Buffers de un mensaje \texttt{WidevineCencHeader} con un
esquema conocido. La llamada \texttt{generateRequest} pasa la caja
PSSH completa al CDM, que parsea el SystemID, los KIDs, y el
header codificado en protobuf para construir un request de
licencia.

Este es el input estructurado atacante-controlable que ingiere la
ranura \texttt{CreateSessionAndGenerateRequest} (ranura 3): un
manifest DASH malicioso o una aplicación EME maliciosa puede
colocar bytes arbitrarios dentro del PSSH y pasarlos al CDM con
control total tanto sobre los campos del header de la caja como
del payload protobuf. Las vulnerabilidades de parser en el camino
PSSH del CDM son alcanzables por atacantes desde cualquier CDN
comprometido que sirve el manifest DASH, desde cualquier JavaScript
comprometido que conduce la aplicación EME, o desde cualquier
atacante MITM que pueda reescribir el manifest DASH en tránsito.

Los dos esquemas estándar de cifrado CENC que el CDM debe soportar
son \textbf{CENC-CTR} (AES-128 en modo Counter, con descriptores
de subsample que indican qué rangos de bytes del track cifrado
están protegidos) y \textbf{CBCS} (AES-128 en modo CBC con
cifrado por \emph{patrón}: un patrón 1-de-10 o 1-de-9 de bloques
cifrados-salteados diseñado para tolerar los límites de NALU de
AVC/HEVC). La estructura \texttt{InputBuffer\_2} de CDM\_11
transporta ambos esquemas y sus parámetros de subsample/patrón
por llamada.

\subsection{Mojo IPC y el proceso utility del CDM}
\label{sec:bg-mojo}

Chromium aísla al CDM en un proceso utility dedicado para
confinar el impacto del compromiso del CDM. El proceso utility se
spawnea con la línea de comando:

\begin{lstlisting}[basicstyle=\ttfamily\tiny]
chrome.exe --type=utility
           --utility-sub-type=media.mojom.CdmServiceBroker
           --lang=en-US --service-sandbox-type=cdm
           --no-sandbox  (when CIG is disabled for research)
           ...
\end{lstlisting}

Dentro del utility, una implementación de interfaz Mojo acepta
los requests CDM del anfitrión y los reenvía a una clase C++ que
posee un único puntero a
\texttt{ContentDecryptionModule\_11}. Ese puntero apunta al objeto
devuelto por
\texttt{widevinecdm.dll!CreateCdmInstance(11, \ldots)} ---
que es el objeto cuya vtable mapeamos en la
Sección~\ref{sec:static}.

\subsection{Control Flow Guard y Code Integrity Guard}
\label{sec:bg-cfg-cig}

Dos mitigaciones a nivel de proceso de Windows son relevantes
para la instrumentación en runtime del utility del CDM:

\textbf{Code Integrity Guard (CIG)} es la política de mitigación
de proceso \texttt{PROCESS\_CREATION\_MITIGATION\_POLICY\_BLOCK\_NON\_\\MICROSOFT\_BINARIES\_ALWAYS\_ON}.
Cuando está habilitada, el kernel rechaza cargar cualquier DLL al
proceso que no esté firmado por Microsoft o un partner de
Microsoft. Chrome habilita CIG en sus procesos renderer y utility
por default desde Chrome 70~\cite{chromium-cig}. La flag
\texttt{--disable-features=RendererCodeIntegrity}, en combinación
con \texttt{--no-sandbox} en los procesos afectados, deshabilita
CIG para uso de investigación.

\textbf{Control Flow Guard (CFG)} es una mitigación a nivel de
compilador de Microsoft que agrega un bitmap validado por kernel
de targets válidos de indirect-call al proceso corriendo e
inserta un check antes de cada indirect call~\cite{microsoft-cfg}.
Una vtable parcheada que apunta a una función hook que recién
asignamos no está en el bitmap por default, así que la primera
llamada por la ranura parcheada dispara un fallo CFG que termina
el proceso.

La forma soportada de marcar una nueva función como target válido
es la API \texttt{SetProcessValidCallTargets}
\cite{ms-spvct-docs}. La API toma un handle de proceso, una base
de página alineada a 4~KB, un tamaño de página, y un array de
descriptores \texttt{CFG\_CALL\_TARGET\_INFO} cuyo campo
\texttt{Offset} es el \emph{offset intra-página} (12 bits bajos)
del target. Una versión temprana de nuestro injector usó los 16
bits bajos, lo cual está mal: la página que respalda al hook es
de 4~KB y el bitmap se indexa con granularidad de 16 bytes
dentro de la página. El injector liberado usa la máscara
correcta:

\begin{lstlisting}[language=C, basicstyle=\ttfamily\scriptsize]
CFG_CALL_TARGET_INFO cti;
cti.Offset = (DWORD)(hookVA & 0xfff);   /* in-page offset */
cti.Flags  = CFG_CALL_TARGET_VALID;
DWORD_PTR pageBase = hookVA & ~0xfffULL;
SetProcessValidCallTargets(hProc, (PVOID)pageBase,
                            0x1000, 1, &cti);
\end{lstlisting}

\subsection{ASLR y la vista compartida del DLL}
\label{sec:bg-aslr}

Un detalle práctico de implementación relevante a nuestro
experimento entre pestañas (Sección~\ref{sec:findings-cross}):
los procesos utility de Chrome cargan
\texttt{widevinecdm.dll} desde la misma imagen en disco cada vez,
y la compartición de image-sections de Windows significa que dos
procesos utility separados ambos terminan con el DLL mapeado en
la \emph{misma dirección virtual}. Por eso las direcciones de
hook de nuestro injector (en la imagen del shim, en la base de
carga reservada del shim) fueron idénticas entre el proceso
utility del CDM de Netflix y el de Bitmovin. Sin embargo, la
sección de datos del DLL es copy-on-write por proceso; cada
utility tiene sus propios contadores y su propio file handle de
log abierto.

% =============================================================================
\section{Modelo de amenaza y metodología}
\label{sec:threat}

\subsection{Modelo de amenaza}
\label{sec:threat-model}

Consideramos a un usuario honesto de un navegador de la familia
Chromium reproduciendo contenido Widevine-protegido con
decodificación por hardware en Windows 11. El usuario ha
consentido a la instrumentación nivel-investigación de su propio
navegador y ha instalado nuestras herramientas de ingeniería
inversa en su propia máquina de laboratorio. El ejecutable del
navegador no está modificado y está firmado por el vendor. El CDM
Widevine es el \texttt{widevinecdm.dll} provisto por el vendor de
la versión del navegador instalado. El usuario posee una
suscripción válida y paga al servicio comercial de streaming
usado en las secciones empíricas del paper, y el uso de esa
suscripción es para visualización personal durante la ventana de
investigación. No hay segundo usuario, ni jurisdicción fuera de la
propia del usuario, ni contenido de terceros involucrado.

Excluimos deliberadamente de este estudio cualquier modelo de
amenaza en el que un atacante capture cuadros descifrados.
Nuestra preocupación no es el potencial de explotación del CDM
sino su forma arquitectural y la superficie de ataque que
presenta al navegador anfitrión y al atacante de red posicionado
entre el navegador y el CDN del streaming.

\subsection{Invariantes de la metodología}
\label{sec:threat-invariants}

Seguimos lo que describimos como el \emph{modelo Buchanan} de
investigación de DRM, resumido anteriormente
(Sección~\ref{sec:ante-buchanan}). La versión operativa del
modelo que gobernó cada decisión experimental en este paper es:

\begin{enumerate}[itemsep=2pt,topsep=2pt]
  \item \textbf{Observación shape-only.} A cualquier hook runtime
        se le permite registrar metadata de despacho --- número de
        ranura, timestamp, hashes de valores de registros, campos
        de tamaño de estructura --- pero se le prohíbe copiar
        cualquier byte apuntado por puntero perteneciente a
        ciphertext, cuadros descifrados, material de clave de
        contenido, o vectores de inicialización, hacia cualquier
        almacenamiento persistente o cualquier canal fuera del
        proceso. El código fuente del shim codifica esta
        invariante explícitamente: no existe camino de código
        que llame a \texttt{memcpy(\&log\_buf, ib.data, \dots)};
        el extractor de forma sólo lee los campos \texttt{size},
        un hash FNV-1a-32 de los primeros 16 bytes de
        \texttt{key\_id}, una flag de presencia distinta-de-cero
        para el IV, y un puñado de campos escalares públicos.
  \item \textbf{Ausencia demostrable.} Después de cada experimento
        verificamos forensicamente que el almacenamiento producido
        es JSON puro ASCII consistente sólo de los campos de
        metadata permitidos y no contiene blobs binarios, ni
        contenido decodificado, ni material de clave. Documentamos
        el procedimiento de verificación en la
        Sección~\ref{sec:ethics} y proveemos el script verificador
        en la metodología liberada.
  \item \textbf{Mutación reversible.} Todas las modificaciones
        runtime al proceso target (parches de vtable,
        instalaciones de hook, registros de bitmap CFG) son
        reversibles. Guardamos el estado pre-parche en formato
        wire al tooling liberado y restauramos el target a su
        estado original antes de publicar. El helper liberado
        retiene las direcciones originales de los thunks de la
        outer-vtable en el array \texttt{g\_pOrigSlot[19]} del
        shim para que un observador externo pueda validar que el
        estado post-experimento es byte-igual al estado
        pre-experimento.
  \item \textbf{Sin superficie de circumvención de contenido.}
        Usamos la propia sesión de navegador del operador a un
        servicio comercial de streaming para confirmar qué hace
        el CDM durante steady-state, pero nunca instrumentamos al
        decoder, ni al componente PMP del lado kernel
        (\texttt{peauth.sys}), ni la superficie protegida, ni
        cualquier otro componente posicionado para materializar
        un cuadro descifrado. El shim shape-only es el hook más
        profundo en el stack que instalamos.
  \item \textbf{Falsificación antes de la publicación.} Cada
        afirmación empírica se somete a una prueba de
        falsificación: una hipótesis alternativa que, si fuera
        verdadera, invalidaría la afirmación, se construye y se
        prueba con cualquier evidencia que la metodología pueda
        producir. La Sección~\ref{sec:robustness} reporta el
        sprint de falsificación.
\end{enumerate}

Una consecuencia en lenguaje natural de estas invariantes es que
nuestros resultados caracterizan el \emph{comportamiento de
despacho} del CDM, no su \emph{estado criptográfico}. Podemos
reportar con qué frecuencia dispara una ranura durante reproducción
real, qué distribución de tuplas de registros ve, cuál es la forma
de su descriptor de buffer entrante. No podemos y no reportamos
ningún byte de contenido decodificado.

% =============================================================================
\section{Ingeniería inversa estática de CDM\_11}
\label{sec:static}

Esta sección recorre la ingeniería inversa estática del despacho
CDM\_11 en \texttt{widevinecdm.dll}, en el orden en que la
realizamos. Dos binarios fueron analizados en paralelo:
\textbf{Chrome 149.0.7827.115}, desde su path de instalación
estándar bajo \texttt{WidevineCdm/\_platform\_specific/win\_x64/}
($\sim$22~MB), y la misma biblioteca tal como viene en
\textbf{Microsoft Edge 150} ($\sim$19~MB). Todos los RVAs se
miden respecto al image base del binario.

\subsection{Paso 1: \texttt{CreateCdmInstance} y el switch de versión}
\label{sec:static-export}

La tabla de exports de \texttt{widevinecdm.dll} contiene seis
entradas; la única relevante para nuestros propósitos es
\texttt{CreateCdmInstance}:

\begin{lstlisting}[basicstyle=\ttfamily\scriptsize]
Edge 150     CreateCdmInstance @ rva 0x1d0850
Chrome 149   CreateCdmInstance @ rva 0x1d2e10
\end{lstlisting}

Disassemblando cualquiera de las dos entradas muestra una cadena
de pequeñas funciones wrapper que cargan una palabra de versión
en \texttt{ecx} (versión 9, 10, u 11 de la interfaz CDM) y llaman
a un constructor por-versión. El disassembly del fragmento
relevante del \texttt{CreateCdmInstance} de Edge 150 es:

\begin{lstlisting}[basicstyle=\ttfamily\scriptsize]
0x1801d0850: push   r14
0x1801d0852: push   rsi
0x1801d0853: push   rdi
0x1801d0854: sub    rsp, 0x40
0x1801d0858: mov    edi, ecx           ; ecx = requested version
0x1801d085a: lea    rax, [rip + 0x...] ; load module-relative base
...
0x1801d0994: mov    ecx, 0xb           ; CDM_11 selected
0x1801d0999: mov    rax, r9            ; arg setup
...
\end{lstlisting}

Escanear ingenuamente el primer \texttt{call} después de cada
\texttt{mov ecx, 0xb} devuelve el stack-check probe del runtime
(\texttt{\_\_chkstk\_ms} o equivalente), no el constructor real.
Un discoverer robusto debe filtrar helpers conocidos del runtime.
El script liberado \texttt{cdm\_arch\_discover.py} usa una
heurística estructural: enumera todos los targets de call
alcanzables desde \texttt{CreateCdmInstance}, filtra targets
llamados más de dos veces (el constructor se llama exactamente
una vez por branch de versión; los helpers de runtime se llaman
muchas veces), y valida cada candidato inspeccionando si su
cuerpo contiene una secuencia de la forma:

\begin{lstlisting}[basicstyle=\ttfamily\scriptsize]
  lea reg, [rip + N]      ; load a candidate vtable RVA
  mov qword ptr [r?], reg ; install at offset 0 of `this`
\end{lstlisting}

\noindent y si el RVA cargado apunta dentro de \texttt{.rdata} a
una ubicación cuyos primeros diecinueve QWORDs son todos punteros
de código válidos dentro de \texttt{.text}. El primer candidato
que cumpla se reporta como el constructor del CDM y el RVA
cargado como su outer vtable. El discoverer localiza el RVA de
outer vtable en Edge 150 en \texttt{0x780250} por este camino; el
outer vtable de Chrome 149 en \texttt{0x768fa0}.

\subsection{Paso 2: Versiones co-empaquetadas del CDM}
\label{sec:static-versions}

Un escaneo \emph{exhaustivo} posterior de \texttt{.rdata} en
busca de estructuras de diecinueve ranuras que coincidan con el
patrón thunk del CDM (descrito abajo) revela \textbf{dos} de esas
estructuras por binario, espaciadas a \texttt{0xc0} bytes:

\begin{center}\small
\begin{tabular}{lcc}
\toprule
\textbf{Binary} & \textbf{vtable A} & \textbf{vtable B} \\
\midrule
Edge 150    & \texttt{0x780190} & \texttt{0x780250} \\
Chrome 149  & \texttt{0x768ee0} & \texttt{0x768fa0} \\
\bottomrule
\end{tabular}
\end{center}

Dentro de cada binario, ambas vtables comparten su entrada de
descifrado por-cuadro (la ranura 14 apunta a la misma función
tanto en A como en B), pero sus entradas \texttt{Initialize}
difieren:

\begin{center}\small
\begin{tabular}{lcc}
\toprule
\textbf{Binary} & \textbf{vtable A.slot[0]} & \textbf{vtable B.slot[0]} \\
\midrule
Edge 150    & 0x1db330 & 0x1db390 \\
Chrome 149  & 0x1e03c0 & 0x1e0420 \\
\bottomrule
\end{tabular}
\end{center}

Esto es consistente con que el binario empaqueta tanto vtables
CDM\_10 como CDM\_11 en soporte a navegadores anfitriones que aún
no se han actualizado a la interfaz más reciente, y coincide con
el historial del header de Chromium. En lo que sigue
identificamos a la \emph{vtable B} como la outer vtable de
CDM\_11, dado que parches en runtime en las posiciones de ranuras
de \emph{vtable B} interceptan el despacho frente al anfitrión
EME (Sección~\ref{sec:runtime}). La implementación compartida en
la ranura 14 refleja la decisión de implementación de usar una
rutina de descifrado por-cuadro a lo largo de ambas versiones de
interfaz; la ranura 0 divergente refleja la firma
\texttt{Initialize} específica de versión.

\subsection{Paso 3: El patrón thunk de la vtable externa}
\label{sec:static-thunk}

La outer vtable contiene diecinueve entradas. La mayoría de ellas
son \emph{thunks} de la forma:

\begin{lstlisting}[basicstyle=\ttfamily\scriptsize]
slot N thunk @ rva 0x1dad00 + (N-3)*0x20:
  mov rcx, qword ptr [rcx + 8]       ; this->m_impl
  mov rax, qword ptr [rcx]           ; m_impl->m_vtable
  mov rax, qword ptr [rax + 8*N]     ; slot N of m_impl_vtable
  mov rdx, qword ptr [rip + cfg_disp]; CFG dispatcher pointer
  jmp rdx                             ; CFG-protected tail-call
\end{lstlisting}

Cada thunk son veintiocho bytes seguidos de relleno
\texttt{int3}, dispuestos en una secuencia uniforme con
espaciado \texttt{0x20} dentro de \texttt{.text} para las ranuras
3 a 17. El thunk de la ranura 14 en Edge 150, disassemblado por
completo:

\begin{lstlisting}[basicstyle=\ttfamily\scriptsize]
0x1801dae60: mov    rcx, qword ptr [rcx + 8]
0x1801dae64: mov    rax, qword ptr [rcx]
0x1801dae67: mov    rax, qword ptr [rax + 0x70]
0x1801dae6b: mov    rdx, qword ptr [rip + 0x10dcd2e]
0x1801dae72: jmp    rdx
0x1801dae74: int3
0x1801dae75: int3
... (int3 padding to 0x20 boundary)
\end{lstlisting}

Disassemblamos cada thunk en ambos binarios y confirmamos el
chequeo de índice de ranura (el offset en \texttt{[rax + 8*N]}
es igual al número de ranura multiplicado por ocho) para quince
de quince ranuras thunkeadas en cada binario. La confirmación
byte-a-byte es lo que licencia nuestra afirmación de
``inactividad de la ranura~14'' en la
Sección~\ref{sec:findings-dormancy}: la evidencia en runtime
($g\_callsBySlot[14] = 0$) es atribuible al método correcto sólo
porque el índice de ranura en el thunk de la outer-vtable
despacha al offset interno correcto.

Las ranuras 0 (\texttt{Initialize}), 1 (\texttt{GetStatusForPolicy})
y 18 (\texttt{OnStorageId}) \emph{no} son thunks; son
implementaciones directas ubicadas en \texttt{.text} en
direcciones no contiguas. La ranura 2
(\texttt{SetServerCertificate}) es un thunk en ambos binarios
pero despacha vía \texttt{[rax + 8]} en lugar de
\texttt{[rax + 0x10]} --- comparte una implementación con la
ranura 1 a nivel de la capa \texttt{m\_impl}. Este es un detalle
arquitectural por-binario consistente tanto en Edge 150 como en
Chrome 149, sugiriendo que es una decisión deliberada de
implementación más que un accidente del build.

\subsection{Paso 4: La vtable \texttt{m\_impl}}
\label{sec:static-mimpl}

El thunk externo despacha hacia una vtable \emph{intermedia}
poseída por el objeto \texttt{m\_impl} que la clase exterior
mantiene en \texttt{this->m\_impl} (\texttt{[outer + 8]}).
Localizar esta vtable intermedia en \texttt{.rdata} es el segundo
paso del análisis estático.

Usamos un pivote de firma conocida. De una traza de call-chain
previa (documentada en nuestras notas de investigación), sabíamos
que el wrapper interno de la ranura 5 (\texttt{UpdateSession}) en
Edge 150 está en RVA \texttt{0x1d4ab0}, identificable por su
prólogo multi-push con un stack canary que carga datos
atacante-controlables a través de una copia camino al parser
ofuscado. El código del wrapper es inconfundible en un listing de
disassembly.

Para encontrar la vtable \texttt{m\_impl}, hacemos una búsqueda
reversa en \texttt{.rdata} del valor QWORD
\texttt{image\_base + 0x1d4ab0}. Hay
exactamente una ocurrencia. Retroceder cinco entradas de ocho
bytes desde el match da la dirección de la vtable
\texttt{m\_impl}: RVA \texttt{0x77d370} en Edge 150. Los
diecinueve QWORDs circundantes apuntan dentro de \texttt{.text},
confirmando al candidato. El mismo pivote aplicado a Chrome 149
(donde el wrapper interno de la ranura 5 está en el RVA análogo)
localiza la vtable \texttt{m\_impl} de Chrome 149 en el offset
correspondiente.

\subsection{Paso 5: La estructura de la capa \texttt{m\_impl}}
\label{sec:static-layer2}

Disassemblando cada entrada de la vtable \texttt{m\_impl} se
revela una mezcla de tres tipos de cuerpo:

\begin{enumerate}[itemsep=2pt,topsep=2pt]
  \item \textbf{Wrappers MSVC limpios} (ranuras 2, 3, 4, 5, 6, 7,
        18). Tienen un prólogo estándar (múltiples
        \texttt{push}~callee-saves, \texttt{sub rsp, N}, carga del
        stack canary), allocan o piden prestadas copias de input,
        realizan saneamiento, y luego llaman a un handler más
        profundo. Estas son las ranuras que ingieren datos de
        longitud variable influenciados por el atacante (strings
        de session-id, payloads de respuesta, buffers de
        init-data). El wrapper de la ranura 3 en Edge 150,
        abreviado para el paper:

\begin{lstlisting}[basicstyle=\ttfamily\scriptsize]
0x1801d47d0: push   r15
0x1801d47d2: push   r14
0x1801d47d4: push   r13
0x1801d47d6: push   r12
0x1801d47d8: push   rsi
0x1801d47d9: push   rdi
0x1801d47da: push   rbx
0x1801d47db: sub    rsp, 0x50
0x1801d47df: mov    esi, r9d            ; init_data_type
0x1801d47e2: mov    edi, r8d            ; session_type
0x1801d47e5: mov    ebx, edx            ; promise_id
0x1801d47e7: mov    eax, [rsp + 0xb8]   ; init_data_size
0x1801d47ee: mov    rdx, [rip + 0x10f284b]
0x1801d47f5: xor    rdx, rsp            ; stack canary
0x1801d47f8: mov    [rsp + 0x48], rdx
0x1801d47fd: mov    r14, [rcx + 0x10]   ; load m_impl-internal
0x1801d4801: xorps  xmm0, xmm0
0x1801d4804: movaps [rsp + 0x30], xmm0
...
0x1801d4821: mov    rcx, r12
0x1801d4824: call   0x180736080         ; malloc/__chkstk
0x1801d4829: mov    [rsp + 0x30], rax
0x1801d482e: lea    r13, [rax + r12]    ; end pointer
...
0x1801d485f: call   0x18021f200         ; LAYER-3 ENTRY
                                        ; (obfuscated parser)
\end{lstlisting}

  \item \textbf{Double-trampolines} (ranuras 0, 1, 8, 9, 11--17).
        Cuerpos de dos instrucciones de la forma
        \texttt{mov rcx, [rcx + N]} seguido por \texttt{jmp <addr>}
        que reenvían la llamada a un handler de capa más profunda
        sin procesamiento intermedio. La entrada de la ranura 14
        de \texttt{m\_impl} en Edge 150:

\begin{lstlisting}[basicstyle=\ttfamily\scriptsize]
0x1801d4f10: mov    rcx, qword ptr [rcx + 0x10]
0x1801d4f14: jmp    0x180230c20    ; layer-3 handler (obfuscated)
\end{lstlisting}

  \item \textbf{Stubs} (ranura 10). Un cuerpo de dos instrucciones
        \texttt{mov eax, 3 ; ret} que devuelve un código de error
        constante sin más trabajo. La ranura 10 es
        \texttt{InitializeAudioDecoder}; en el build L1 de Windows
        es intencionalmente un no-op porque la decodificación de
        audio es manejada por el pipeline Media Foundation fuera
        del CDM.
\end{enumerate}

\noindent La vtable \texttt{m\_impl}, en otras palabras, es
mayormente una capa de \emph{ruteo}; el trabajo real sucede una
indirección más adelante.

\subsection{Paso 6: El modelo de despacho de tres capas}
\label{sec:static-3layer}

Componiendo los cinco pasos anteriores, la cadena de despacho
desde una llamada virtual del anfitrión hasta un handler
por-ranura se muestra en la Figura~\ref{fig:3layer}. Cada llamada
virtual pública atraviesa tres indirecciones antes de alcanzar la
función que realiza el trabajo.

\begin{figure}[t]
\centering
\resizebox{\columnwidth}{!}{%
\begin{tikzpicture}[node distance=0.55cm]
  \node[cdmbox, fill=gray!10, minimum width=10cm] (host) {Chromium media/cdm host (renderer)};
  \node[layerbox, fill=blue!10, below=of host] (l1) {%
    \textbf{Layer 1: \texttt{outer\_vtable} @ \texttt{.rdata + 0x780250}}\\[2pt]
    \scriptsize 17 thunks $+$ 3 direct impls.
    Thunk: \texttt{mov rcx,[rcx+8]; mov rax,[rcx]; mov rax,[rax+8N]; jmp [CFG]}};
  \node[layerbox, fill=green!10, below=of l1] (l2) {%
    \textbf{Layer 2: \texttt{m\_impl} vtable @ \texttt{.rdata + 0x77d370}}\\[2pt]
    \scriptsize 7 clean MSVC wrappers (slots 2--7, 18) $+$
    11 double-trampolines (slots 0, 1, 8, 9, 11--17) $+$ 1 stub (slot 10).};
  \node[layerbox, fill=red!12, below=of l2] (l3) {%
    \textbf{Layer 3: per-slot handlers (scattered across \texttt{.text})}\\[2pt]
    \scriptsize Where parsing, cryptographic processing, and state-machine work happen.
    MBA-style obfuscation concentrated on slots that parse attacker-controllable input.};
  \draw[fatarrow] (host.south) -- node[right,font=\scriptsize,xshift=2pt] {virtual call via CDM\_11 interface} (l1.north);
  \draw[fatarrow] (l1.south) -- node[right,font=\scriptsize,xshift=2pt] {slot $N \cdot 8$ dispatch} (l2.north);
  \draw[fatarrow] (l2.south) -- node[right,font=\scriptsize,xshift=2pt] {trampoline target / wrapper call} (l3.north);
\end{tikzpicture}%
}%
\caption{El modelo de despacho de tres capas de CDM\_11. Cada
llamada virtual pública del CDM desde el anfitrión atraviesa tres
indirecciones antes de alcanzar la función que realiza el
trabajo. Los RVAs mostrados son de Edge 150
\texttt{widevinecdm.dll}; Chrome 149 tiene la misma estructura en
RVAs correspondientemente diferentes (ver
Sección~\ref{sec:cross}).}
\label{fig:3layer}
\end{figure}

\subsection{Una nota sobre los falsos arranques}
\label{sec:static-false}

Para el beneficio de investigadores que sigan este trabajo, el
análisis estático tal como en realidad lo llevamos a cabo
incluyó dos falsos arranques substanciales.

El primer falso arranque fue identificar una tabla de despacho
singleton interna en \texttt{.data + 0x1588b48} como la
``outer vtable de CDM\_11''. Esa tabla, de hecho, es una
superficie de despacho interna separada de diecinueve ranuras
dentro de la biblioteca Widevine que no es la interfaz pública
EME. Nuestra conclusión inicial de que ``la ranura por-cuadro de
la outer vtable de CDM\_11 dispara $\sim$600 veces por segundo
bajo reproducción HW de Netflix en steady-state''
era incorrecta: era la ranura 0 del singleton interno la que
disparaba 600 veces por segundo, y la ranura por-cuadro real de
la outer-vtable estaba en cero. La corrección vino de un pivote
de traza de constructor que localizó la outer vtable real vía el
export \texttt{CreateCdmInstance}, tras lo cual el hallazgo de
inactividad pasó a ser correctamente atribuible a la superficie
de cara al usuario.

El segundo falso arranque fue la atribución original de
ofuscación estilo MBA a la capa \texttt{m\_impl}. La capa
\texttt{m\_impl} de hecho es una capa de ruteo limpia (según la
Sección~\ref{sec:static-layer2}); la ofuscación vive en la capa
tres. Una enumeración capa-por-capa de las diecinueve ranuras,
documentada en nuestras notas de investigación 13b y 13c,
localizó la ofuscación en su profundidad real.

Destacamos estos falsos arranques porque son típicos del trabajo
de RE estático sobre binarios cerrados con indirecciones
internas de despacho. Una lectura en una sola pasada del call
graph del binario naturalmente se habría asentado o en la vtable
incorrecta o en la capa incorrecta, y el sprint de falsificación
de la Sección~\ref{sec:robustness} es lo que atrapó ambos.

% =============================================================================
\section{Instrumentación en runtime}
\label{sec:runtime}

El mapa estático de la sección anterior arroja hipótesis sobre
qué ranuras son alcanzadas durante reproducción en steady-state y
cuáles no. Para verificar esas hipótesis instrumentamos el
proceso utility del CDM en ejecución con un hook shape-only.

\subsection{Code Integrity Guard (CIG) y el flag de launch}
\label{sec:runtime-cig}

Chrome habilita Code Integrity Guard
(\texttt{PROCESS\_CREATION\_MITIGATION\_POLICY\_BLOCK\_NON\_\\MICROSOFT\_BINARIES\_ALWAYS\_ON})
en sus procesos utility desde Chrome 70~\cite{chromium-cig}. CIG
rechaza cualquier \texttt{LoadLibrary} de una imagen no firmada
por Microsoft a nivel del kernel, por lo que una inyección
vainilla \texttt{CreateRemoteThread(LoadLibraryA)} de nuestro
shim falla en la etapa \texttt{LdrLoadDll} con
\texttt{STATUS\_DYNAMIC\_CODE\_BLOCKED} (0xC0000605).

Para uso de investigación en laboratorio, Chrome puede ser
lanzado con las flags
\texttt{--disable-features=}\linebreak[1]\texttt{RendererCodeIntegrity}
y \texttt{--no-sandbox}, que en conjunto deshabilitan CIG en los
procesos utility. Esto produce un target de investigación
controlado sin modificar la biblioteca CDM ni el binario del
navegador anfitrión, y corresponde a la asunción equivalente de
alcanzabilidad que un atacante con ejecución de código en el
renderer tendría que satisfacer por separado (vía, por ejemplo,
una arbitrary-write al JIT del renderer o una cadena de
process-load que evada CIG).

\subsection{Registro del bitmap de Control Flow Guard}
\label{sec:runtime-cfg}

CIG-deshabilitado no deshabilita Control Flow Guard (CFG): una
vez que el shim está cargado, simplemente parchar una entrada de
vtable para que apunte a una función hook no registrada causa
que el dispatcher CFG aborte la llamada como una violación de
control de flujo. El installer liberado usa
\texttt{SetProcessValidCallTargets} para marcar cada hook target
como un call target válido antes de parchar su ranura. La
Figura~\ref{fig:hookinstall} bosqueja el flujo de instalación.

\begin{figure}[t]
\centering
\begin{tikzpicture}[node distance=0.55cm and 0.7cm, font=\small]
  \node[cdmbox, fill=blue!10, minimum width=4.5cm] (s1) {\textbf{1.} Launch Chrome with\\\scriptsize \texttt{--disable-features=}\\\scriptsize \texttt{RendererCodeIntegrity --no-sandbox}};
  \node[cdmbox, fill=blue!10, below=of s1, minimum width=4.5cm] (s2) {\textbf{2.} CDM utility process spawns,\\\scriptsize loads \texttt{widevinecdm.dll}};
  \node[cdmbox, fill=orange!15, below=of s2, minimum width=4.5cm] (s3) {\textbf{3.} \texttt{CreateRemoteThread(LoadLibraryA)}\\\scriptsize injects \texttt{wv\_shim.dll} into utility};
  \node[cdmbox, fill=red!15, below=of s3, minimum width=4.5cm] (s4) {\textbf{4.} For each slot N to hook:\\\scriptsize Read \texttt{outer\_vtable[N]} \(\rightarrow\) save to \texttt{g\_pOrigSlot[N]}};
  \node[cdmbox, fill=red!15, below=of s4, minimum width=4.5cm] (s5) {\textbf{5.} \texttt{SetProcessValidCallTargets(}\\\scriptsize \texttt{ page\_base, page\_size, hook\_va \& 0xfff, VALID)}\\\scriptsize Register \texttt{WvHookSlotN} as CFG-valid};
  \node[cdmbox, fill=green!15, below=of s5, minimum width=4.5cm] (s6) {\textbf{6.} \texttt{VirtualProtectEx + WriteProcessMemory}\\\scriptsize Patch \texttt{outer\_vtable[N] = WvHookSlotN}};
  \node[cdmbox, fill=green!15, below=of s6, minimum width=4.5cm] (s7) {\textbf{7.} Hook fires on next CDM virtual call.\\\scriptsize Shim logs shape, calls original via \texttt{g\_pOrigSlot[N]}};
  \foreach \a/\b in {s1/s2, s2/s3, s3/s4, s4/s5, s5/s6, s6/s7}
    \draw[fatarrow] (\a) -- (\b);
\end{tikzpicture}
\caption{Secuencia de instalación del hook. El Paso~5 (registro
del bitmap CFG) es lo que atrapa a los investigadores con más
frecuencia --- el campo \texttt{Offset} de
\texttt{CFG\_CALL\_TARGET\_INFO} espera un offset intra-página
(12 bits bajos) del hook target, no una dirección absoluta ni
otra máscara. El Paso~4 captura la sutileza de diseño que
motivó al array de back-pointers por-ranura
(\S\ref{sec:runtime-orig}).}
\label{fig:hookinstall}
\end{figure}

\subsection{Diseño del shim shape-only}
\label{sec:runtime-shape}

El shim, \texttt{wv\_shim.dll}, exporta diecinueve funciones hook
por-ranura \texttt{WvHookSlot0} hasta \texttt{WvHookSlot18}, cada
una con su propio slot de puntero original en un array exportado
\texttt{g\_pOrigSlot[19]}, y un único array contador
\texttt{g\_callsBySlot[19]}. Cada función hook:

\begin{enumerate}[itemsep=2pt,topsep=2pt]
  \item Atomically increments its slot counter and the aggregate
        \texttt{g\_calls}.
  \item Calls the original via \texttt{g\_pOrigSlot[N]} with the
        unmodified register tuple.
  \item For slots 9, 14, 15 (which take an \texttt{InputBuffer\_2*}
        argument in \texttt{rdx}), tries to read the buffer's
        \emph{shape} via the safe-copy primitive described below
        and writes a single JSONL line with
        \texttt{encryption\_scheme}, \texttt{data\_size},
        \texttt{key\_id\_size}, \texttt{iv\_size}, an FNV-1a-32 hash
        of the first 16 bytes of \texttt{key\_id}, an
        \texttt{iv\_has\_nonzero} bit, \texttt{num\_subsamples},
        \texttt{pattern\_crypt}, \texttt{pattern\_skip}, and the
        media \texttt{timestamp}.
  \item For other slots, writes a \texttt{generic} event with an
        FNV-1a-32 hash of the four argument registers.
\end{enumerate}

\noindent No byte of \texttt{ib.data} (the ciphertext) is ever read,
either transiently or persistently; \texttt{ib.key\_id} and
\texttt{ib.iv} are read into a stack-local buffer that is collapsed
into a hash and then immediately discarded. The shim source enforces
this by simply not having a code path that copies any of those
fields into a logging structure.

El primitivo safe-copy usa \texttt{VirtualQuery} para confirmar
que cada página fuente está commit'eada y es legible antes de
copiar, ya que el compilador C de \texttt{mingw-w64} no provee
el SEH \texttt{\_\_try / \_\_except} de Microsoft en x86\_64. El
check recorre cada página del rango fuente, confirma el estado
\texttt{MEM\_COMMIT}, y confirma una de las protecciones de
página con permiso de lectura (\texttt{PAGE\_READONLY},
\texttt{PAGE\_READWRITE}, \texttt{PAGE\_EXECUTE\_READ}, etc.)
antes del \texttt{memcpy}.

\subsection{Cadena de punteros originales por ranura}
\label{sec:runtime-orig}

Una sutileza de diseño que vale la pena notar: la primera versión
del shim usaba una sola variable \texttt{g\_pOrigDecrypt} para
todas las ranuras hookeadas. Parchar múltiples ranuras con ese
único back-pointer llevaba a crashes por signature-mismatch en el
momento que disparaba la segunda ranura parcheada (el back-pointer
contenía el original de la ranura anterior). El shim liberado usa
un array por-ranura. Aún más sutilmente, cuando hookeamos la
vtable \emph{exterior} en lugar de una tabla singleton interna,
el back-pointer debe contener la dirección del thunk exterior,
\emph{no} una dirección de tabla de despacho interna; si no, el
camino de llamada-original termina saltando a la función interna
incorrecta y corrompe silenciosamente la llamada. Nuestro
installer (\texttt{wv\_helper3.c}) lee el original de la outer
vtable en el momento del parche y lo guarda en
\texttt{g\_pOrigSlot[N]} como parte de la misma operación
atómica:

\begin{lstlisting}[language=C, basicstyle=\ttfamily\scriptsize]
DWORD_PTR origFn = 0;
ReadProcessMemory(hp, (LPCVOID)slotAddr, &origFn, 8, NULL);

/* CFG: register hook as valid target */
CFG_CALL_TARGET_INFO cti = { hookVA & 0xfff,
                              CFG_CALL_TARGET_VALID };
SetProcessValidCallTargets(hp, (PVOID)(hookVA & ~0xfffULL),
                            0x1000, 1, &cti);

/* Save outer thunk address into g_pOrigSlot[N] */
WriteProcessMemory(hp,
    (LPVOID)(gOrigSlotRemote + N*8),
    &origFn, 8, NULL);

/* Patch outer_vtable[N] -> hook */
DWORD oldProt = 0;
VirtualProtectEx(hp, (LPVOID)slotAddr, 8,
                  PAGE_READWRITE, &oldProt);
WriteProcessMemory(hp, (LPVOID)slotAddr, &hookVA, 8, NULL);
VirtualProtectEx(hp, (LPVOID)slotAddr, 8, oldProt, &oldProt);
\end{lstlisting}

% =============================================================================
\section{Hallazgos empíricos}
\label{sec:findings}

Reportamos cuatro hallazgos, cada uno anclado a una corrida
experimental grabada. Todas las corridas usaron Chrome
149.0.7827.115 en Windows 11 26100 con decoder NVIDIA
hardware-secure disponible. Ningún proceso utility del CDM
crasheó durante las corridas que produjeron los datos reportados
aquí.

\subsection{Inactividad bajo decode por hardware en la superficie pública}
\label{sec:findings-dormancy}

\textbf{Hipótesis.} Bajo reproducción L1 con decode por hardware,
las ranuras de descifrado por-cuadro de la outer vtable de
CDM\_11 (ranuras 9, 14, 15) no son invocadas.

\textbf{Procedimiento.} Chrome fue lanzado con
\texttt{--disable-features=RendererCodeIntegrity --no-sandbox}.
El operador se autenticó en Netflix y comenzó un título 4K HDR
que la plataforma reporta como decoded por hardware. Después de
noventa segundos de reproducción en steady-state instalamos
nuestro shim en el proceso utility del CDM y parchamos cuatro
ranuras de la outer-vtable simultáneamente: la ranura 14
(\texttt{DecryptAndDecodeFrame}), la 15
(\texttt{DecryptAndDecodeSamples}), la 3
(\texttt{CreateSessionAndGenerateRequest}, control), y el RVA
del thunk exterior de la ranura 14. Luego observamos los
contadores correspondientes por noventa segundos más de
reproducción continua.

\textbf{Resultado.} Los cuatro contadores se mantuvieron en
cero. El video continuó reproduciéndose normalmente todo el
tiempo. La inactividad de la ranura 14 se mantuvo incluso a
través de recuperaciones (restauramos la ranura 14 a su original
a mitad del experimento y
re-patched; the counter remained at zero across both windows). The
four-slot zero-fires-in-ninety-seconds result was reproducible
3/3.

\textbf{Interpretación.} Durante reproducción Netflix
hardware-decoded en steady-state, la superficie por-cuadro de la
interfaz CDM\_11 no se ejerce. Un atacante que comprometa el
renderer anfitrión y se posicione sobre la vtable virtual de cara
a EME no puede interceptar cuadros decodificados a través de la
interfaz CDM\_11, porque los cuadros no pasan por ella.

\begin{figure}[t]
\centering
\resizebox{\columnwidth}{!}{%
\begin{tikzpicture}[node distance=0.45cm and 0.8cm]
  \node[cdmbox, fill=purple!12, minimum width=4cm] (cdn) {Encrypted DASH/CMAF stream};
  \node[cdmbox, fill=gray!10, below left=0.7cm and 0.3cm of cdn, minimum width=4cm] (sw_route) {SW decode path};
  \node[cdmbox, fill=gray!10, below right=0.7cm and 0.3cm of cdn, minimum width=4cm] (hw_route) {HW decode path (L1)};
  \node[cdmbox, fill=red!15, below=of sw_route, minimum width=4cm] (sw_cdm) {CDM\_11 \texttt{Decrypt} / \texttt{DecryptAndDecodeFrame}\\\scriptsize (would be invoked here)};
  \node[cdmbox, fill=blue!10, below=of sw_cdm, minimum width=4cm] (sw_render) {Renderer in browser process};
  \node[cdmbox, fill=red!10, dashed, below=of hw_route, minimum width=4cm] (hw_cdm) {CDM\_11 \texttt{DecryptAndDecodeFrame}\\\scriptsize \textbf{DORMANT} (\(g\_{}calls{=}0\))};
  \node[cdmbox, fill=orange!15, below=of hw_cdm, minimum width=4cm] (mfpmp) {Media Foundation PMP\\\scriptsize \texttt{peauth.sys} / hardware decoder};
  \node[cdmbox, fill=cyan!15, below=of mfpmp, minimum width=4cm] (gpu) {Protected GPU surface};
  \draw[thinarrow] (cdn) -- (sw_route);
  \draw[thinarrow] (cdn) -- (hw_route);
  \draw[thinarrow] (sw_route) -- (sw_cdm);
  \draw[thinarrow] (sw_cdm) -- (sw_render);
  \draw[thinarrow] (hw_route) -- (hw_cdm);
  \draw[thinarrow, dashed] (hw_cdm) -- node[right,font=\scriptsize] {bypass} (mfpmp);
  \draw[thinarrow] (mfpmp) -- (gpu);
  \draw[fatarrow, red, dashed] (hw_route.east) .. controls +(2,-1) and +(2,1) .. (mfpmp.east)
    node[midway, right, font=\scriptsize, text=red, align=left] {per-frame\\bypass};
\end{tikzpicture}%
}%
\caption{Inactividad bajo decode por hardware. Bajo decode SW
(camino izquierdo), las ranuras por-cuadro de CDM\_11 son el
target obvio. Bajo decode HW L1 (camino derecho), el ciphertext
por-cuadro es ruteado a través de Media Foundation PMP y el
decoder de kernel, materializándose directamente sobre una
superficie GPU protegida; las ranuras por-cuadro de CDM\_11
permanecen dormidas.}
\label{fig:hwdormancy}
\end{figure}

\subsection{Actividad del singleton interno a $\sim$2.3k ops/s}
\label{sec:findings-singleton}

\textbf{Hipótesis.} A pesar de que la superficie CDM\_11 esté
dormida por-cuadro, la biblioteca CDM sí realiza trabajo durante
reproducción en steady-state. Ese trabajo es alcanzable a través
de una tabla de despacho interna separada, distinta de la outer
vtable.

\textbf{Procedimiento.} Parchamos las ranuras 0--7 de la tabla
singleton interna identificada en \texttt{.data + 0x1588b48} en
Chrome 149, luego observamos durante $\sim$355 segundos de
reproducción Netflix continua.

\textbf{Resultado.} El shim grabó \textbf{810{,}099 eventos} en un
único log JSONL (en dos capturas consecutivas). Distribución:

\begin{center}\small
\begin{tabular}{rrrr}
\toprule
\textbf{slot} & \textbf{count} & \textbf{unique \texttt{rh}} & \textbf{rate (ev/s)} \\
\midrule
0 & 395{,}580 & 16{,}338 & $\sim$1{,}114 \\
2 &     20    &  2     & 0.06 \\
6 & 414{,}459 & 1{,}379 & $\sim$1{,}167 \\
\bottomrule
\end{tabular}
\end{center}

\noindent La tasa agregada es $\sim$2{,}281 eventos/seg en esta
tabla interna durante reproducción HW en steady-state.

\begin{figure}[h]
\centering
\begin{tikzpicture}
\begin{axis}[
  width=0.95\columnwidth,
  height=4.2cm,
  ybar,
  bar width=10pt,
  ylabel={Calls per second},
  ylabel style={font=\small},
  xlabel={Internal singleton slot index},
  xlabel style={font=\small},
  symbolic x coords={0,1,2,3,4,5,6,7},
  xtick=data,
  ymin=0, ymax=1300,
  ymajorgrids,
  grid style=dashed,
  nodes near coords,
  nodes near coords style={font=\tiny, /pgf/number format/.cd, fixed, precision=0},
  enlarge x limits=0.08,
  font=\footnotesize
]
\addplot[fill=red!30] coordinates {(0,1114) (1,0) (2,0.06) (3,0) (4,0) (5,0) (6,1167) (7,0)};
\end{axis}
\end{tikzpicture}
\caption{Perfil de actividad de la tabla de despacho singleton
interna en \texttt{.data + 0x1588b48} durante reproducción
Netflix 4K HDR en steady-state.}
\label{fig:singleton-rate}
\end{figure}

\textbf{Interpretación.} El despacho singleton interno
\emph{no} es la interfaz CDM\_11 de cara al usuario; la tasa de
disparo de la ranura 0 es incompatible con la semántica de
\texttt{Initialize} (que sería llamada una sola vez por sesión,
no 600 veces por segundo). Lo más plausible es que la tabla
represente despacho estilo worker OEMCrypto, un anillo de
allocator, o una escalera de primitivas criptográficas interna
a la biblioteca Widevine. De todas formas, esto es invisible
desde el límite EME y desde el despacho de la outer-vtable.
Combinado con el hallazgo de inactividad anterior, el cuadro es:
\emph{la biblioteca CDM está haciendo trabajo criptográfico
constante durante reproducción HW, pero ese trabajo está
escondido detrás de una superficie de despacho diferente a la que
conoce el cableado EME de Chrome}.

\subsection{Alcanzabilidad de las ranuras de gestión de sesión}
\label{sec:findings-session}

\textbf{Hipótesis.} A pesar de que las ranuras por-cuadro de la
outer-vtable estén dormidas bajo decode HW, las ranuras de
gestión-de-sesión de la outer-vtable
(\texttt{CreateSessionAndGenerateRequest} y
\texttt{UpdateSession}) siguen siendo alcanzables y \emph{sí}
disparan durante eventos normales de ciclo de vida de sesión.

\textbf{Procedimiento.} Parchamos las ranuras 3 y 5 de la
outer-vtable con back-pointers por-ranura correctos a sus
direcciones originales de thunk
(Sección~\ref{sec:runtime-orig}). El operador realizó un cambio
de título en la pestaña Netflix en ejecución y observamos los
contadores \texttt{g\_callsBySlot} durante el minuto siguiente.

\textbf{Resultado.}

\begin{center}\small
\begin{tabular}{lcc}
\toprule
& \texttt{slot 3} & \texttt{slot 5} \\
\midrule
pre-action & 10 & 9 \\
post-action & \textbf{13 (+3)} & \textbf{12 (+3)} \\
\bottomrule
\end{tabular}
\end{center}

\noindent Tres invocaciones de CreateSessionAndGenerateRequest y
tres de UpdateSession fueron observadas en los segundos siguientes
al cambio de título. El proceso utility del CDM se mantuvo vivo
($\sim$85~MB working set), como se esperaba dado que el
back-pointer contenía la dirección correcta del thunk.

\textbf{Interpretación.} Las ranuras de gestión-de-sesión de la
outer vtable no están dormidas; se ejercen durante transiciones
de sesión iniciadas por el usuario y se mantienen como una
superficie de ataque real para input estructurado
atacante-controlable (cajas PSSH vía el manifest DASH; payloads
de respuesta-de-licencia vía el CDN de licencias).

\subsection{Consistencia entre fuentes (Bitmovin)}
\label{sec:findings-cross}

\textbf{Hipótesis.} El modelo de despacho y la alcanzabilidad de
ranuras son propiedades del CDM Widevine, no de Netflix como una
aplicación EME particular.

\textbf{Procedimiento.} Después del cambio de título de Netflix
anterior, el operador abrió una pestaña separada en el mismo
navegador Chrome y cargó el demo público Widevine de Bitmovin en
\texttt{bitmovin.com/demos/drm}. Chrome spawnó un \emph{segundo}
proceso utility CDM para servir esa pestaña. Instalamos nuestro
shim en ese segundo proceso y parchamos las ranuras 3 y 5 de su
outer-vtable. El operador recargó el demo para forzar la
re-creación de sesión.

\textbf{Resultado.} Contadores del segundo proceso utility CDM:

\begin{center}\small
\begin{tabular}{lcc}
\toprule
& \texttt{slot 3} & \texttt{slot 5} \\
\midrule
post-reload & \textbf{2} & \textbf{2} \\
\bottomrule
\end{tabular}
\end{center}

\begin{figure}[t]
\centering
\resizebox{\columnwidth}{!}{%
\begin{tikzpicture}[node distance=0.4cm and 0.5cm, font=\small]
  \node[cdmbox, fill=gray!10, minimum width=10cm] (browser) {Single Chrome browser process (two tabs)};
  \node[cdmbox, fill=blue!10, below left=0.5cm and -0.5cm of browser, minimum width=4cm] (nf) {Tab 1: Netflix\\\scriptsize (4K HDR HW playback)};
  \node[cdmbox, fill=cyan!12, below right=0.5cm and -0.5cm of browser, minimum width=4cm] (bm) {Tab 2: Bitmovin demo\\\scriptsize (public Widevine sample)};
  \node[cdmbox, fill=orange!15, below=of nf, minimum width=4cm] (u1) {CDM utility \texttt{pid=5772}\\\scriptsize \texttt{widevinecdm.dll @ 0x7ffc...}};
  \node[cdmbox, fill=orange!15, below=of bm, minimum width=4cm] (u2) {CDM utility \texttt{pid=5048}\\\scriptsize \texttt{widevinecdm.dll @ 0x7ffc...}};
  \node[cdmbox, fill=green!15, below=of u1, minimum width=4cm] (h1) {Outer vtable patched:\\\scriptsize slot 3 \(\rightarrow\) WvHookSlot3\\\scriptsize slot 5 \(\rightarrow\) WvHookSlot5};
  \node[cdmbox, fill=green!15, below=of u2, minimum width=4cm] (h2) {Outer vtable patched:\\\scriptsize slot 3 \(\rightarrow\) WvHookSlot3\\\scriptsize slot 5 \(\rightarrow\) WvHookSlot5};
  \node[cdmbox, fill=yellow!20, below=of h1, minimum width=4cm] (r1) {After title switch:\\\scriptsize \textbf{slot 3: +3 fires}\\\scriptsize \textbf{slot 5: +3 fires}};
  \node[cdmbox, fill=yellow!20, below=of h2, minimum width=4cm] (r2) {After F5 + play:\\\scriptsize \textbf{slot 3: +2 fires}\\\scriptsize \textbf{slot 5: +2 fires}};
  \draw[thinarrow] (browser.south) -- (nf.north);
  \draw[thinarrow] (browser.south) -- (bm.north);
  \draw[thinarrow] (nf) -- (u1);
  \draw[thinarrow] (bm) -- (u2);
  \draw[thinarrow] (u1) -- (h1);
  \draw[thinarrow] (u2) -- (h2);
  \draw[thinarrow] (h1) -- (r1);
  \draw[thinarrow] (h2) -- (r2);
\end{tikzpicture}%
}%
\caption{Experimento entre fuentes. Dos pestañas del mismo
navegador Chrome reproducen dos fuentes Widevine-protegidas
distintas, cada una servida por su propio proceso utility CDM.
Hooks idénticos confirman alcanzabilidad idéntica de las ranuras
de gestión-de-sesión.}
\label{fig:crosssource}
\end{figure}

\noindent Dos invocaciones de CreateSessionAndGenerateRequest y
dos de UpdateSession fueron observadas. El modelo de despacho
(thunk de la outer-vtable, wrapper m\_impl, handler de capa-3) es
idéntico a lo que observamos en el proceso utility del CDM que
sirve a Netflix; los caminos de código de los hooks de la ranura
3 y la ranura 5 fueron ejercidos en la misma secuencia en ambos
procesos.

\textbf{Interpretación.} El modelo de despacho de CDM\_11 es una
propiedad del binario \texttt{widevinecdm.dll}, no una propiedad
de la aplicación EME que lo usa. Un atacante posicionado para
influenciar el payload PSSH o el de licencia de \emph{cualquier}
fuente Widevine-protegida verá el mismo despacho y la misma
alcanzabilidad de ranura que uno posicionado contra Netflix.

\subsection{Observación arquitectural lateral: utilities CDM separados}

Una observación secundaria que vale la pena grabar es que Chrome
149 spawnó \emph{dos} procesos utility CDM cuando Netflix y
Bitmovin estaban reproduciéndose simultáneamente en pestañas
diferentes, no un utility compartido. Esto es consistente con la
política de aislamiento per-origin de Chromium e implica que un
atacante posicionado para comprometer el CDM de una aplicación
EME no puede observar directamente el CDM de una segunda
aplicación EME vía estado in-process compartido.

% =============================================================================
\section{La distribución de ofuscación como capa defensiva}
\label{sec:obfuscation}

El binario \texttt{widevinecdm.dll} contiene numerosas funciones
cuyo cuerpo coincide con la firma de libro de texto de la
ofuscación mixed-boolean-arithmetic (MBA): constantes mágicas
grandes residentes en stack escritas inmediatamente al entrar;
instrucciones de avance de estado por hash multiplicativo de la
forma \texttt{imul edx, ecx, 0xK ; shr edx, 0xS ; cmp eax, 0xT
; je state\_exit}; cascadas de despacho sin branches con
\texttt{cmovae};
patrones de relleno \texttt{0xaaaaaaaa} o
\texttt{0xaaaaaaaaaaaaaaaa} actuando como valores de
inicialización poison para estado local del stack. Estos
patrones no están uniformemente distribuidos a lo largo del
espacio de ranuras.

\subsection{Disposición visual}

La Figura~\ref{fig:obfuscation-grid} visualiza el estado de
ofuscación por-ranura lado a lado con la clasificación de cada
ranura como ranura de input atacante o ranura passthrough.

\begin{figure}[t]
\centering
\resizebox{\columnwidth}{!}{%
\begin{tikzpicture}[%
  slotcell/.style={draw, minimum width=0.62cm, minimum height=0.55cm, font=\scriptsize, align=center, inner sep=1.5pt},
  legendcell/.style={draw, minimum width=0.42cm, minimum height=0.42cm, inner sep=0pt},
]
  % Legend row at top
  \begin{scope}[shift={(0, 1.55)}]
    \node[font=\footnotesize, anchor=west] at (-7.0, 0) {Per-slot CDM\_11 obfuscation status:};
    \node[legendcell, fill=red!30, anchor=west] at (-2.0, 0) {};
    \node[font=\scriptsize, anchor=west] at (-1.4, 0) {obfuscated};
    \node[legendcell, fill=orange!30, anchor=west] at (0.2, 0) {};
    \node[font=\scriptsize, anchor=west] at (0.8, 0) {partial};
    \node[legendcell, fill=green!25, anchor=west] at (2.0, 0) {};
    \node[font=\scriptsize, anchor=west] at (2.6, 0) {clean};
    \node[legendcell, fill=gray!20, anchor=west] at (3.8, 0) {};
    \node[font=\scriptsize, anchor=west] at (4.4, 0) {tramp/stub};
  \end{scope}

  % Header letters (A/P) above the cells
  \foreach \i/\inp in {
    0/A, 1/A, 2/P, 3/A, 4/A, 5/A, 6/-, 7/-, 8/-, 9/A,
    10/-, 11/-, 12/-, 13/-, 14/A, 15/A, 16/A, 17/-, 18/P}
  {
    \pgfmathsetmacro\xpos{\i*0.7 - 6.3}
    \ifx\inpA
      \node[font=\tiny] at (\xpos, 0.55) {\textbf{A}};
    \else\ifx\inpP
      \node[font=\tiny] at (\xpos, 0.55) {\textbf{P}};
    \fi\fi
  }

  % Slot cells - single row, properly spaced
  \foreach \i/\col in {
    0/orange!30, 1/red!30, 2/green!25, 3/red!30, 4/gray!20,
    5/red!30, 6/gray!20, 7/gray!20, 8/red!30, 9/red!30,
    10/gray!20, 11/orange!30, 12/green!25, 13/green!25, 14/red!30,
    15/green!25, 16/green!25, 17/green!25, 18/green!25}
  {
    \pgfmathsetmacro\xpos{\i*0.7 - 6.3}
    \node[slotcell, fill=\col] at (\xpos, 0) {\textbf{\i}};
  }

  % Method names below, smaller and clearly separated
  \foreach \i/\name in {
    0/Init, 1/GetStat, 2/SetCert, 3/CreateSess, 4/LoadSess,
    5/UpdateSess, 6/CloseSess, 7/RemoveSess, 8/TimerExp, 9/Decrypt,
    10/InitAud, 11/InitVid, 12/DeinitDec, 13/ResetDec, 14/DecDecFr,
    15/DecDecSm, 16/PlatChal, 17/QryOutP, 18/StoreId}
  {
    \pgfmathsetmacro\xpos{\i*0.7 - 6.3}
    \node[font=\tiny, rotate=60, anchor=east] at (\xpos, -0.45) {\name};
  }

  % Footnote line at bottom
  \node[font=\scriptsize, align=center, text width=12cm] at (0, -2.5) {%
    \textbf{A} = attacker-input slot (parses structured data);\quad
    \textbf{P} = passthrough (opaque blob the CDM forwards without parsing).};
\end{tikzpicture}%
}%
\caption{Estado de ofuscación por-ranura de CDM\_11. Las celdas
rojas son ranuras cuyo cuerpo de capa-1 o capa-3 coincide con la
firma de ofuscación estilo MBA; las naranjas son matches
parciales; las verdes son limpias a lo largo de la cadena de
despacho; las grises son trampolines de capa-2 o stubs cuyos
handlers de capa-3 reportamos individualmente en la
Sección~\ref{sec:obfuscation}.}
\label{fig:obfuscation-grid}
\end{figure}

\subsection{El patrón}

Las ranuras que obviamente pasan un \emph{blob} opaco
influenciado por el atacante que el CDM reenvía sin parsear
(\texttt{SetServerCertificate}, \texttt{OnStorageId}) están
limpias en cada capa del despacho. Las ranuras que comienzan a
parsear datos estructurados atacante-controlables exhiben
ofuscación, aunque a profundidades variables. La ranura que toma
una política textual (\texttt{GetStatusForPolicy}) está ofuscada
en la capa 1, donde comienza el parsing. Las ranuras que toman
payloads estructurados grandes
(\texttt{CreateSessionAndGenerateRequest},
\texttt{UpdateSession}) colocan el procesamiento ofuscado uno o
dos saltos más profundo, después de que un wrapper MSVC limpio
ha copiado el input al buffer propio del CDM y buscado la
sesión. Las ranuras involucradas en decode por-cuadro
(\texttt{Decrypt}, \texttt{DecryptAndDecodeFrame}) están
ofuscadas inmediatamente en el target del double-trampoline,
donde comienza el procesador criptográfico.

Este es el patrón que cuidadosamente resumimos como:

\begin{quote}\itshape
CDM\_11 protege cada ranura a la profundidad en la que esa
ranura comienza a parsear datos estructurados del atacante.
\end{quote}

\noindent Es una observación arquitectural, no una afirmación de
seguridad: la ofuscación no es seguridad, y no estamos evaluando
la fuerza de la protección, sólo su distribución.

% =============================================================================
\section{Estabilidad entre binarios}
\label{sec:cross}

El modelo arquitectural descrito en la
Sección~\ref{sec:static-3layer} fue desarrollado sobre Edge 150
\texttt{widevinecdm.dll}. Validamos que el mismo modelo se
mantiene en el \texttt{widevinecdm.dll} de Chrome 149 corriendo
nuevamente el discoverer:

\begin{center}\small
\begin{tabular}{lcc}
\toprule
\textbf{Property} & \textbf{Edge 150} & \textbf{Chrome 149} \\
\midrule
CreateCdmInstance RVA          & 0x1d0850 & 0x1d2e10 \\
CDM\_11 outer vtable RVA       & 0x780250 & 0x768fa0 \\
CDM\_10 outer vtable RVA       & 0x780190 & 0x768ee0 \\
Slot-index check (slots 3--17) & 15/15    & 15/15    \\
Direct-impl slots              & 3 (0, 1, 18) & 3 (0, 1, 18) \\
\bottomrule
\end{tabular}
\end{center}

\noindent El modelo también generaliza a través de versiones de
interfaz del CDM dentro de un único binario: tanto en Edge 150
como en Chrome 149, las vtables CDM\_10 y CDM\_11 exhiben un
despacho idéntico de tres capas. Su handler de descifrado
por-cuadro es la misma función física; sus handlers
\texttt{Initialize} divergen. Esto es consistente con que el
binario provea ambas versiones de interfaz para
compatibilidad-hacia-atrás con navegadores anfitriones más
viejos, y con que el autor del CDM eligió compartir la
implementación de los métodos estables-en-versión entre las dos.

La tabla de ranuras inter-binarios aparece como el
Apéndice~\ref{app:slot-table}.

% =============================================================================
\section{Pruebas de robustez}
\label{sec:robustness}

Antes de publicar cualquier afirmación hecha más arriba,
corrimos un sprint de falsificación contra los hallazgos
empíricos más fuertes. El sprint incluyó:

\begin{enumerate}[itemsep=2pt,topsep=2pt]
  \item \textbf{Corrección del índice de ranura.} El thunk de la
        ranura 14 de la outer-vtable está byte-confirmado para
        despachar en offset \texttt{0x70} ($= 14 \cdot 8$) en la
        vtable \texttt{m\_impl}.
  \item \textbf{Señal vs ruido en los datos de forma.} La captura
        de 810k eventos fue analizada por entropía del hash de
        registros y distribución temporal. La entropía del hash
        para las dos ranuras calientes fue 16{,}338 y 1{,}379
        valores únicos a lo largo de 395{,}580 y 414{,}459
        invocaciones respectivamente; la tasa fue estable a lo
        largo de la ventana de observación a $\sim$2.3k
        eventos/s.
  \item \textbf{Enumeración de capa-3.} Los diecinueve handlers
        de capa-3 fueron localizados vía el target del jump del
        trampoline o el primer call (no-stack-check) del wrapper
        y fueron clasificados por estado de ofuscación. La
        primera versión de la afirmación de ofuscación, que
        ubicaba la ofuscación en la capa \texttt{m\_impl}, fue
        encontrada incorrecta; la ofuscación está un salto más
        profundo, y la afirmación fue reescrita en consecuencia.
  \item \textbf{Identidad inter-binarios.} El discoverer fue
        re-corrido contra el \texttt{widevinecdm.dll} de Chrome
        149. La outer vtable, el mapeo índice-de-ranura, y la
        identificación de método CDM\_11 por-ranura coincidieron
        con el análisis de Edge 150.
  \item \textbf{Identidad entre fuentes.} La instrumentación de
        la outer-vtable fue desplegada contra el segundo proceso
        utility del CDM spawnado para una fuente Widevine
        no-Netflix (Bitmovin). Las ranuras 3 y 5 dispararon en
        secuencia idéntica al utility de Netflix.
  \item \textbf{Conteo de instancias CDM.} La única pestaña de
        Netflix spawnó exactamente un proceso utility CDM a lo
        largo de una ventana de observación de ocho horas; no
        observamos un segundo utility materializándose para la
        misma pestaña.
\end{enumerate}

\subsection{Hallazgos candidatos refutados}
\label{sec:robustness-refuted}

El sprint de falsificación también sacó a la superficie dos
hallazgos candidatos que fueron refutados por análisis byte-level
más profundo. Los documentamos aquí porque tipifican los
candidatos de clase-de-bug que una metodología menos rigurosa
hubiera falsamente promovido:

\textbf{P1 (escritura sin autenticación al pool de PEAuth):}
Una auditoría de subagente de \texttt{peauth.sys} inicialmente
marcó una aparente escritura no autenticada al pool del kernel
en offset \texttt{0x140042980} donde un campo de tamaño
controlado por el usuario en \texttt{[rdi+0x14]} parecía ser
usado como length para \texttt{memcpy} antes del chequeo de
firmado \texttt{CiGetPEInformation} en \texttt{0x1400429ae}. El
disassembly byte-level del rango \texttt{0x42900..0x429e0}
reveló un chequeo de overflow de 32-bit explícito
(\texttt{lea r12d, [r15+rsi]; cmp r12d, r15d; jb fail}) en
\texttt{0x42964} y un chequeo de cota superior en
\texttt{0x4296d} entre la lectura del parámetro y el memcpy. El
candidato es refutado; no hay bug del lado kernel presente.

\textbf{W2 (lectura OOB en el parser PSSH):} Una auditoría de
subagente de CDM\_11 marcó los offsets \texttt{0x1D6E1},
\texttt{0x1D8CF}, \texttt{0x1DB9D}, y \texttt{0x1DFB} como
candidatos para una lectura out-of-bounds en el parsing del
tamaño del header PSSH/MP4-box
(\texttt{mov r10d, [r{si/di}+4]}). El disassembly en esos
offsets reveló no parsing de cajas sino la round-table para una
transformación MD5, con los operandos siendo XORs con constantes
de estado y no lecturas de buffer. El candidato es refutado; los
bytes son código cripto, no código de parser.

Estas refutaciones y un candidato adicional que se mantuvo
inconcluyente establecieron el principio metodológico de que
cualquier afirmación de subagente ``byte-verified de alta
confianza'' debe pasar una revisión manual byte-a-byte
independiente antes de ser promovida a un hallazgo grado-de-
publicación.

% =============================================================================
\section{Manual para el defensor}
\label{sec:defender}

Esta sección traduce los hallazgos arquitecturales y empíricos
en acciones defensivas concretas. El lector objetivo es un ingeniero
de seguridad protegiendo un deployment de navegadores de la familia
Chromium, un incident responder triando un compromiso sospechado
del CDM, o un ingeniero de detección construyendo hooks para un EDR.

\subsection{Tabla de superficie de ataque por ranura}

Table~\ref{tab:attack-surface} summarises which of the nineteen
CDM\_11 outer-vtable slots are an attack surface in practice,
which are dormant during normal hardware-decoded playback, and
which are unreachable from a network-positioned attacker.

\begin{table}[h]
\centering\small
\begin{tabular}{rlll}
\toprule
\textbf{Ranura} & \textbf{Método} & \textbf{Alcance atacante} & \textbf{Actividad en steady-state} \\
\midrule
0  & Initialize                       & host & once per session \\
1  & GetStatusForPolicy               & host & rare \\
2  & SetServerCertificate             & host-cert & rare \\
3  & CreateSessionAndGenerateRequest  & \textbf{manifest} & once per session start \\
4  & LoadSession                      & host & rare \\
5  & UpdateSession                    & \textbf{license CDN} & on license rotate \\
6  & CloseSession                     & host & on teardown \\
7  & RemoveSession                    & host & rare \\
8  & TimerExpired                     & internal & internal \\
9  & Decrypt                          & per-frame  & dormant (HW path) \\
10 & InitializeAudioDecoder           & host & stub returns 3 \\
11 & InitializeVideoDecoder           & host & once \\
12 & DeinitializeDecoder              & host & on teardown \\
13 & ResetDecoder                     & host & rare \\
14 & DecryptAndDecodeFrame            & per-frame  & dormant (HW path) \\
15 & DecryptAndDecodeSamples          & per-frame  & dormant (HW path) \\
16 & OnPlatformChallengeResponse      & host & rare \\
17 & OnQueryOutputProtectionStatus    & host & rare \\
18 & OnStorageId                      & host & rare \\
\bottomrule
\end{tabular}
\caption{Superficie de ataque por-ranura y perfil de actividad
en steady-state. ``Alcance atacante'' es la fuente controlable
más fuerte de los argumentos de la ranura:
\textbf{manifest} = caja PSSH DASH/CMAF atacante-controlable
desde el CDN que sirve el manifest; \textbf{license CDN} =
respuesta del servidor de licencias atacante-controlable (MITM
en el CDN o servicio de licencia comprometido);
\textbf{per-frame} = ciphertext desde los segmentos de medio del
stream; \emph{host} = valor opaco provisto por el navegador
anfitrión; \emph{internal} = scheduling interno del CDM.
``Actividad en steady-state'' resume si la ranura dispara durante
reproducción en curso una vez que la sesión está establecida.}
\label{tab:attack-surface}
\end{table}

\subsection{Firmas de detección}

Tres firmas de detección salen directamente de la arquitectura
y de la evidencia en runtime.

\textbf{Firma D1 --- canario de vtable.} La outer vtable vive en
la sección de sólo-lectura \texttt{.rdata} de
\texttt{widevinecdm.dll}. Cada una de las diecinueve entradas
debería apuntar dentro de \texttt{.text} de
\texttt{widevinecdm.dll}. Un check de integridad in-process en
un EDR que lea los diecinueve QWORDs de la outer-vtable después
de la inicialización del CDM y verifique que todos resuelven
dentro del intervalo
\texttt{[image\_base, image\_base+image\_size)} de
\texttt{widevinecdm.dll} atrapará el patrón de manipulación en
runtime usado por el shim de este paper y por cualquier otra
herramienta de hooking del CDM que use la misma forma.

\textbf{Firma D2 --- drift del bitmap CFG.} Cada llamada a
\texttt{SetProcessValidCallTargets} dentro de un proceso utility
CDM es en sí una señal. En operación normal, el utility CDM no
auto-registra targets CFG-válidos adicionales; el bitmap se
establece en module-load y no crece. Un evento de crecimiento
después del steady-state del proceso, particularmente uno cuyos
nuevos targets válidos están fuera del rango de imagen de
\texttt{widevinecdm.dll}, es un indicador fuerte de un intento
de inyección-y-parche.

\textbf{Firma D3 --- actividad en ranuras por-cuadro en camino
HW.} Nuestro hallazgo empírico
(Sección~\ref{sec:findings-dormancy}) es que las ranuras 9, 14, y
15 están dormidas durante reproducción HW-decoded en
steady-state. Una sonda in-process que hookee esas tres ranuras
y observe cualquier actividad durante reproducción HW ha
atrapado o bien a un atacante forzando fallback a decode SW
(una técnica conocida de CDN hostiles) o a un atacante in-renderer
que aún no se dio cuenta que el camino por-cuadro está dormido.

\subsection{Flujo de triage ante sospecha de CDM comprometido}

\begin{enumerate}[itemsep=2pt]
  \item Identificar el proceso utility CDM sospechado vía
        \texttt{tasklist /m widevinecdm.dll}; en una sesión de
        reproducción única-fuente Widevine debería haber uno.
  \item Verificar el RVA de la outer-vtable contra el binario
        usando el discoverer liberado
        (\S\ref{sec:static-export}). Hacer cross-check de la
        outer vtable cargada en runtime para tampering usando
        la Firma D1.
  \item Para cada una de las ranuras 3, 5, 14, 15, leer el
        QWORD en la entrada de la outer-vtable y confirmar que
        es el thunk original. Las ranuras 14 y 15 en particular
        son los caminos por-cuadro a los que un atacante
        in-renderer apuntaría si no estuviera al tanto de la
        inactividad bajo decode HW.
  \item Examinar el bitmap CFG en busca de adiciones fuera del
        rango de imagen de \texttt{widevinecdm.dll} (Firma D2).
\end{enumerate}

% =============================================================================
\section{Lectura para el atacante}
\label{sec:attacker}

Esta sección traduce los mismos hallazgos desde el otro lado de
la línea. El lector objetivo es un investigador de seguridad mapeando
la superficie de ataque EME, un operador de red-team scopeando un
follow-on a nivel kernel, o un participante de coordinated-disclosure
que necesita saber qué ranuras merecen tiempo de auditoría profunda.

\subsection{Qué es alcanzable y qué no}

El hallazgo de inactividad
(Sección~\ref{sec:findings-dormancy}) es mala noticia para el
ataque obvio: un renderer comprometido no puede interceptar
cuadros decodificados a través de la interfaz CDM\_11 durante el
decode por hardware L1 porque las ranuras por-cuadro no son
invocadas. El atacante que quiere cuadros decodificados tiene
que encontrar una capacidad independiente contra Media Foundation
PMP, contra el allocator de superficie protegida del driver de
GPU, o contra el lado user-mode del contrato OPM que enforza
\texttt{peauth.sys}. Esos son targets de clase
kernel-y-driver, un orden de magnitud más difíciles de alcanzar
desde un compromiso del renderer.

Las ranuras de gestión-de-sesión son buena noticia. La ranura 3
(\texttt{CreateSessionAndGenerateRequest}) ingiere la caja PSSH,
que es atacante-controlable desde cualquier CDN comprometido que
sirva el manifest DASH/CMAF. La ranura 5
(\texttt{UpdateSession}) ingiere la respuesta del servidor de
licencias, que es atacante-controlable desde cualquier endpoint
de licencia comprometido o cualquier atacante MITM-posicionado.
Ambas ranuras despachan a través de un wrapper MSVC limpio en la
capa \texttt{m\_impl} y luego llaman a un handler de capa-3
MBA-ofuscado. El wrapper limpio es auditable; el handler de
capa-3 es el target realista de bug-hunting.

\subsection{Dónde auditar}

Priorización concreta para tiempo de auditoría de capa-3:

\begin{enumerate}[itemsep=2pt]
  \item \textbf{Capa-3 de ranura 3} (Edge 150 RVA
        \texttt{0x21f200}). Parser de PSSH/MP4-box,
        atacante-controlable desde el manifest. Fuertemente
        MBA-ofuscado. La herramienta correcta aquí es fuzzing
        estructurado de mutaciones de la caja PSSH antes que
        disasm manual.
  \item \textbf{Capa-3 de ranura 5} (wrapper limpio en Edge 150
        \texttt{0x1d4ab0}, despacha al procesador whitebox-AES
        ofuscado en \texttt{0x233bd0} vía una cola de
        worker-thread). Parser de respuesta-de-licencia,
        atacante-controlable. Camino async; más difícil de
        alcanzar con fuzzing sincrónico.
  \item \textbf{Capa-3 de ranura 9} (RVA \texttt{0x22ec00}).
        Handler de \texttt{Decrypt}. Fuertemente MBA-ofuscado.
        Según el hallazgo de inactividad, dormido en el camino
        HW; sólo sería alcanzable en un escenario de decode SW.
  \item \textbf{Capa-3 de ranura 14} (RVA \texttt{0x230c20}).
        \texttt{DecryptAndDecodeFrame}. El más fuertemente
        ofuscado de todos los handlers que medimos (puntaje 22
        en nuestro clasificador). Mismo caveat de inactividad que
        la ranura 9.
\end{enumerate}

\subsection{Lo que no es el objetivo correcto}

La ranura 2 (\texttt{SetServerCertificate}) y la ranura 18
(\texttt{OnStorageId}) ambas pasan blobs opacos que el CDM
reenvía sin parsear. La cadena de despacho limpia en ambas capas
es la señal empírica de que no hay nada que auditar ahí. Gastar
tiempo de auditoría en estas ranuras es consistente con el modo
de fallo de auditorías que asumen que todas las ranuras de
camino-crítico están ofuscadas; no lo están.

% =============================================================================
\section{Generalizabilidad}
\label{sec:generalisability}

El CDM Widevine L1 en Windows es uno de tres runtimes comerciales
de DRM de escritorio que un investigador podría ser invitado a
mapear. Microsoft PlayReady viene en el IE-mode de Edge y dentro
del stack
\texttt{Windows.Media.Protection.MediaProtectionManager}. Apple
FairPlay viene en Safari sobre macOS y dentro de TV.app. Ambos
exponen un contrato host-CDM que es, al nivel de abstracción al
que opera este paper, muy similar al contrato CDM\_11 de
Widevine: una cantidad chica de métodos host-callable, una
biblioteca opaca que construye una instancia de una interfaz
versionada, un camino interno de despacho que hace el trabajo
real.

La metodología que liberamos está parametrizada para que
generalice a estos targets con mínima modificación:

\begin{itemize}[itemsep=2pt]
  \item El discoverer \texttt{cdm\_arch\_discover.py} localiza
        el patrón de instalación-de-vtable de un constructor en
        cualquier binario PE por búsqueda estructural; el único
        supuesto Widevine-específico es el nombre del export
        \texttt{CreateCdmInstance}, que un investigador
        aplicando esto a PlayReady o FairPlay sustituye por el
        entry point relevante.
  \item La semántica de hook del shim shape-only es
        independiente de la firma por-ranura; cada
        \texttt{WvHookSlotN} toma la tupla completa de
        callee-save de registros. La misma forma aplica a una
        vtable de PlayReady o FairPlay que usa una forma de
        firma diferente por ranura.
  \item El bypass de CFG vía \texttt{SetProcessValidCallTargets}
        es una mitigación a nivel de proceso de Windows, no a
        nivel Widevine. Aplica sin cambios a cualquier proceso
        de Windows que cargue una biblioteca de tercero y
        habilite CFG.
\end{itemize}

Un paper separado sobre cada uno de PlayReady y FairPlay
aplicando esta metodología es, esperamos, trabajo follow-on que
futuros investigadores producirán. El modelo Buchanan que
restringe la metodología a observación shape-only es el mismo
en los tres casos.

\subsection*{Para investigadores que sigan}

Dos de las decisiones de diseño de este paper están pensadas
para cortocircuitar caminos sin salida comunes en trabajo
follow-on.

Primero, la metodología está estructurada para ser aplicada a
cualquier build futuro de \texttt{widevinecdm.dll} sin edición
manual de RVAs. El discoverer localiza la outer vtable desde el
binario solo; el script de instrumentación en runtime toma el
RVA descubierto como un parámetro. Re-correr la misma metodología
contra un Chrome stable futuro, o contra las variantes
\texttt{libwidevinecdm.so} que vienen en Linux y Chrome OS, es
mayormente cuestión de correr el discoverer.

Segundo, la narrativa de falsos arranques en la
Sección~\ref{sec:static-false} está pensada para evitar dos
caminos sin salida comunes que se ven superficialmente correctos.
Los investigadores entrando a este espacio de RE deberían esperar
encontrarse con varias superficies internas de despacho de
diecinueve ranuras que \emph{no} son la vtable pública de
CDM\_11, y deberían esperar que la capa de
atribución-de-ofuscación del primer intento esté incorrecta en
una indirección. El pivote de traza-de-constructor desde
\texttt{CreateCdmInstance} es el desambiguador confiable.

% =============================================================================
\section{Ética, disclosure, y lo que no hicimos}
\label{sec:ethics}

Este estudio no contiene:
\begin{itemize}[itemsep=1pt,topsep=2pt]
\item ningún cuadro descifrado, en ninguna forma, capturado a
      disco, enviado por IPC, ni retenido transitoriamente a
      través del run del shim;
\item ningún byte de ciphertext, transitorio o persistente,
      recuperado a través del shim;
\item ningún material de clave, clave de contenido, key wrap, ni
      clave de cifrado de contenido, capturado en ninguna capa;
\item ningún ataque contra \texttt{peauth.sys}, el PMP de
      Windows, ni cualquier componente del lado decoder
      posicionado para materializar un cuadro decodificado;
\item ningún código de exploit, shim weaponizado, ni herramienta
      de circumvención de contenido de estante;
\item ningún ataque de red sobre producción contra Netflix,
      Bitmovin, ni ninguna aplicación EME estudiada;
\item ningún resultado de ingeniería inversa que no podamos
      reproducir desde la metodología solo-fuente liberada
      contra el mismo binario.
\end{itemize}

Hemos verificado forensicamente la ausencia de bytes de contenido
en cada artefacto que produjimos. Los logs de forma capturados
son JSON puro ASCII totalizando $\sim$56~MB y contienen sólo los
campos de metadata descritos en la
Sección~\ref{sec:runtime-shape}. El conteo de hashes únicos de
\texttt{key\_id} en los eventos runtime capturados es cero: el
camino de código que hubiera computado un hash de un key\_id de
Netflix nunca fue ejercido en runtime, porque las ranuras que
toman un argumento \texttt{InputBuffer\_2*} no fueron invocadas.

Ningún vendor fue notificado antes de esta submission porque el
trabajo no constituye un reporte de vulnerabilidad. Coordinaremos
con el Chrome Vulnerability Reward Program de Google y los
canales de seguridad de Widevine relevantes antes de cualquier
submission futura de hallazgos clase-exploit derivados de la
metodología aquí.

% =============================================================================
\section{Conclusión}
\label{sec:conclusion}

Hemos presentado un modelo arquitectural verificado byte-a-byte
del despacho moderno de Widevine CDM\_11 tal como viene en dos
navegadores representativos de la familia Chromium, y un mapa
empírico runtime-confirmado de qué ranuras de CDM\_11 se
ejercen durante reproducción Netflix hardware-decoded en
steady-state. Las ranuras por-cuadro de cara al usuario están
dormidas bajo decode por hardware L1 mientras una tabla de
despacho interna genera $\sim$2.3k eventos por segundo de
actividad criptográfica o de scheduler que es invisible desde
el límite EME. Las ranuras de gestión-de-sesión siguen siendo
alcanzables y ejercen el mismo camino de despacho sin importar
si la aplicación EME anfitriona es Netflix o una fuente de
prueba Widevine no-comercial. La biblioteca distribuye
ofuscación mixed-boolean-arithmetic a lo largo del espacio de
ranuras en un patrón que refleja la profundidad a la que cada
ranura comienza a parsear datos estructurados
influenciados-por-atacante.

El trabajo está situado en dos décadas de investigación previa
de DRM, sigue el modelo Buchanan de publicación
sólo-arquitectural, y no produce artefacto de extracción. La
metodología y todo el código fuente necesario para reproducir
los resultados empíricos se liberan junto al paper.

% =============================================================================
\bibliographystyle{plain}
\begin{thebibliography}{99}

\bibitem{widevine-spec}
Google.
\textit{Widevine DRM Architecture Overview.}
Google Developers documentation, accessed 2026.

\bibitem{chromium-cdm-api}
Chromium Project.
\textit{\texttt{media/cdm/api/content\_decryption\_module.h}.}
Chromium source tree, accessed 2026.

\bibitem{chromium-cdm-mojo}
Chromium Project.
\textit{\texttt{media.mojom.ContentDecryptionModule} interface.}
Chromium source tree, accessed 2026.

\bibitem{chromium-cig}
Chromium Project.
\textit{Code Integrity Guard on utility processes.}
Chromium documentation, accessed 2026.

\bibitem{microsoft-cfg}
Microsoft Corporation.
\textit{Control Flow Guard.}
Microsoft Learn documentation, accessed 2026.

\bibitem{ms-spvct-docs}
Microsoft Corporation.
\textit{SetProcessValidCallTargets function.}
Microsoft Learn documentation, accessed 2026.

\bibitem{w3c-eme}
W3C.
\textit{Encrypted Media Extensions, W3C Recommendation 18 September 2017.}

\bibitem{w3c-pssh}
W3C.
\textit{ISO Common Encryption EME Stream Format and Initialization Data.}

\bibitem{iso-cenc}
ISO/IEC 23001-7.
\textit{Information technology — MPEG systems technologies — Part 7: Common encryption in ISO base media file format files.}

\bibitem{buchanan-blog}
D.~Buchanan.
\textit{Personal blog: \url{https://www.da.vidbuchanan.co.uk/}.}
Series of architectural reverse-engineering writeups, 2019--present.

\bibitem{buchanan-imessage}
D.~Buchanan.
\textit{Analysing the iMessage NaCl key-wrap.}
Personal blog, 2023.

\bibitem{buchanan-t2}
D.~Buchanan.
\textit{Apple T2 secure boot chain analysis.}
Personal blog, 2022.

\bibitem{bunnie-xbox-book}
A.~Huang.
\textit{Hacking the Xbox: An Introduction to Reverse Engineering.}
No Starch Press, 2003.

\bibitem{bunnie-nettv}
A.~Huang and S.~Cross.
\textit{NeTV: A Stream Mediation Device, and the DMCA Exemption Process.}
Project documentation, 2011--2015.

\bibitem{sklyarov-case}
\textit{United States v. ElcomSoft Ltd., 203 F. Supp. 2d 1111 (N.D. Cal. 2002).}

\bibitem{dvdjon-norway}
\textit{Det Offentlige v. Jon Lech Johansen.}
Borgarting Court of Appeal, Norway, 2003.

\bibitem{decss-2600}
\textit{Universal City Studios, Inc. v. Reimerdes, 111 F. Supp. 2d 294 (S.D.N.Y. 2000).}

\bibitem{hadad-case}
Reporting and analysis of Israeli prosecutions of DRM reverse engineers
in the early 2000s, secondary sources.

\bibitem{quarkslab-widevine}
Quarkslab security research team.
\textit{Reverse engineering of Widevine: a year of research.}
Quarkslab technical blog, 2019.

\bibitem{quarkslab-android-widevine}
Quarkslab security research team.
\textit{Inside Widevine on Android: the OEMCrypto TrustZone trampoline.}
Quarkslab technical blog, 2020.

\bibitem{l3-decryptor}
Independent contributors.
\textit{Widevine L3 decryptor: reduction of software-only Widevine to a clean-room C implementation.}
GitHub, 2019--present.

\bibitem{ango-l3-decryptor}
Independent contributors.
\textit{ANGO L3 Widevine decryptor.}
GitHub, 2021--present.

\bibitem{checkpoint-widevine}
Check Point Research.
\textit{Various analyses of Widevine on Android.}
Check Point Research blog, 2019--2024.

\bibitem{lief-tool}
R.~Thomas et al.
\textit{LIEF: Library to Instrument Executable Formats.}
Quarkslab, \url{https://lief.re/}, 2017--present.

\bibitem{eyrolles-mba}
N.~Eyrolles, L.~Goubin, M.~Videau.
\textit{Defeating MBA-based Obfuscation.}
Proceedings of the 2016 ACM Workshop on Software PROtection.

\bibitem{biondi-mba}
F.~Biondi, S.~Josse, A.~Legay, T.~Sirvent.
\textit{Effectiveness of synthesis in concolic deobfuscation.}
Computers \& Security, vol. 70, 2017.

\bibitem{chow-aes-whitebox}
S.~Chow, P.~Eisen, H.~Johnson, P.C.~van~Oorschot.
\textit{White-Box Cryptography and an AES Implementation.}
SAC 2002.

\bibitem{lepoint-whitebox}
T.~Lepoint, M.~Rivain.
\textit{White-box cryptography: a survey.}
IACR ePrint, 2014.

\bibitem{capstone-tool}
N.A.~Quynh et al.
\textit{Capstone: Next-generation disassembly framework.}
Black Hat USA 2014.

\bibitem{pefile-tool}
E.~Carrera.
\textit{pefile: Portable Executable format parser.}
GitHub, 2008--present.

\end{thebibliography}

\appendix

% =============================================================================
\section{Tabla de despacho por ranura, ambos binarios}
\label{app:slot-table}

\begin{center}\footnotesize
\begin{tabular}{rlll}
\toprule
\textbf{Slot} & \textbf{Method} & \textbf{Edge 150 RVA} & \textbf{Chrome 149 RVA} \\
\midrule
0  & Initialize                       & 0x1db390 & 0x1e0420 \\
1  & GetStatusForPolicy                & 0x1da7a0 & 0x1df8d0 \\
2  & SetServerCertificate              & 0x1dab10 & 0x1dfba0 \\
3  & CreateSessionAndGenerateRequest   & 0x1dad00 & 0x1dfdc0 \\
4  & LoadSession                       & 0x1dad20 & 0x1dfde0 \\
5  & UpdateSession                     & 0x1dad40 & 0x1dfe00 \\
6  & CloseSession                      & 0x1dad60 & 0x1dfe20 \\
7  & RemoveSession                     & 0x1dad80 & 0x1dfe40 \\
8  & TimerExpired                      & 0x1dada0 & 0x1dfe60 \\
9  & Decrypt                           & 0x1dadc0 & 0x1dfe80 \\
10 & InitializeAudioDecoder            & 0x1dade0 & 0x1dfea0 \\
11 & InitializeVideoDecoder            & 0x1dae00 & 0x1dfec0 \\
12 & DeinitializeDecoder               & 0x1dae20 & 0x1dfee0 \\
13 & ResetDecoder                      & 0x1dae40 & 0x1dff00 \\
14 & DecryptAndDecodeFrame             & 0x1dae60 & 0x1dff20 \\
15 & DecryptAndDecodeSamples           & 0x1dae80 & 0x1dff40 \\
16 & OnPlatformChallengeResponse       & 0x1daea0 & 0x1dff60 \\
17 & OnQueryOutputProtectionStatus     & 0x1daec0 & 0x1dff80 \\
18 & OnStorageId                       & 0x1db310 & 0x1e03a0 \\
\bottomrule
\end{tabular}
\end{center}

% =============================================================================
\section{Esquema del log de forma}
\label{app:schema}

Los eventos JSONL capturados usan una de dos formas.

\textbf{Llamada genérica} (para ranuras que no consumen
\texttt{InputBuffer\_2}):
\begin{lstlisting}[basicstyle=\ttfamily\footnotesize]
{
  "kind" : "call",
  "slot" : <int 0..18>,
  "n"    : <monotonic call number for this slot>,
  "tsc"  : <QueryPerformanceCounter snapshot>,
  "rh"   : <FNV-1a-32 hash of rcx XOR rdx XOR r8 XOR r9>
}
\end{lstlisting}

\textbf{Forma} (para las ranuras 9, 14, 15):
\begin{lstlisting}[basicstyle=\ttfamily\footnotesize]
{
  "kind" : "ib",
  "slot" : <9|14|15>,
  "n"    : <monotonic call number>,
  "tsc"  : <QueryPerformanceCounter snapshot>,
  "es"   : <encryption_scheme; 0=clear, 1=CENC-CTR, 2=CBCS>,
  "is"   : <data_size; size only, never the bytes>,
  "kis"  : <key_id_size>,
  "kh"   : <FNV-1a-32 hash of first 16 bytes of key_id>,
  "ivs"  : <iv_size>,
  "ivnz" : <1 if IV has any non-zero byte, else 0>,
  "ns"   : <num_subsamples>,
  "pc"   : <pattern_crypt>,
  "ps"   : <pattern_skip>,
  "mt"   : <media timestamp>
}
\end{lstlisting}

Campos notablemente ausentes:
\texttt{data}, bytes de \texttt{key\_id}, bytes de \texttt{iv},
array \texttt{subsamples}, contenido del cuadro de salida,
clave de cifrado de contenido, clave de key-wrap, estado interno
de OEMCrypto. Estos están deliberadamente omitidos.

\end{document}
